\chapter{Le graphe de churn et la preuve multi-site}
\label{chap:churn}

\begin{objectifs}
  \item Comprendre pourquoi une boucle de rendu répartie sur plusieurs effets échappe
        au point fixe d'un composant et au bras \og self-churn\fg{} de \regle{infinite-loop}.
  \item Savoir lire et construire le \terme{graphe de churn} : slots qualifiés, arêtes
        \code{Must} / \code{May}, auto-arêtes des effets sans deps, partition \code{self_slot}.
  \item Suivre la recherche de cycles (Tarjan en deux passes) et la stratification
        \Error{} / \Warning{} qu'elle induit.
  \item Maîtriser le \emph{kill de convergence} : de la condition \og un seul écrivain\fg{}
        (\adr{18}) à la preuve multi-site \code{converges_under_all_writes} (\adr{42}), puis
        aux \emph{sites} et aux \emph{revivers} du commit \snapshotCommit.
  \item Comprendre pourquoi l'ensemble des sites convergents est un \emph{plus petit} point
        fixe, et pourquoi le plus grand serait insound.
  \item Situer la frontière de certification : ce que le moteur affirme (\code{Must}/\code{May})
        et ce que la règle re-dérive avant de frapper une \Error.
\end{objectifs}

\prerequis{\cref{chap:relations} (la relation \relation{slot_writers}, ses colonnes
\code{owner}, \code{block}, \code{guard_block}, \code{written} ; la relation
\relation{effect_triggers} ; la preuve mono-site \code{converges_once_written}, ses gardes et
ses trois bras), \cref{chap:statevalue-stabilite} (le treillis \code{Stability},
\code{PerRender} et \code{Versioned}), \cref{chap:point-fixe-composant} (environnements
convergés, tranche finale du point fixe), \cref{chap:interpretation-abstraite} (must / may,
dominance, plus petit point fixe).}

Ce chapitre est le plus difficile de la partie moteur. Il traite d'une question qu'aucun
des mécanismes vus jusqu'ici ne sait trancher seul : \emph{deux} effets, chacun anodin,
peuvent former ensemble une boucle de rendu infinie ; et, symétriquement, un effet qui
\emph{a l'air} de boucler peut converger parce qu'une garde meurt une fois sa propre
écriture posée dans le slot. Le premier problème demande un graphe ; le second une preuve.
Et les deux se rencontrent : une preuve de convergence peut être \emph{ranimée} par une
écriture voisine, si bien que la preuve doit connaître \emph{tous} les sites qui écrivent
le slot, et savoir lesquels ne tirent qu'une fois. Nous suivrons l'ordre historique, qui
est aussi l'ordre de difficulté : le graphe sans kill, le kill mono-écrivain, la preuve
multi-site, les sites, les revivers et le plus petit point fixe.

Tout ce chapitre vit dans deux fichiers du moteur, \file{src/engine/churn.rs}
(927 lignes) et \file{src/engine/guards.rs} (1\,058 lignes), tous deux nés au commit
\code{05d3573} (\adr{42}) et profondément remaniés au commit \snapshotCommit. Les tests
d'intégration sont dans \file{tests/effect_cycles.rs}. Toutes les sorties d'analyseur
montrées ici ont été observées sur des fichiers de \file{/tmp/ch15/} avec
\code{reactant check <fichier> --rule infinite-loop --rule cross-component-infinite-loop --info --show-clean --trace} ;
les tables de relations et d'arêtes proviennent d'une petite sonde qui appelle
\code{analyze_program} puis \code{ProgramRelations::new(&result).churn()} et imprime les
lignes (colonnes élaguées à l'affichage, jamais modifiées).

% ---------------------------------------------------------------------------
\section{Le problème : deux effets qui se relancent}
\label{sec:churn-probleme}

\subsection{Un programme de huit lignes}

\begin{exemple}{Deux effets miroirs}{churn-two-effects}
Chaque effet recopie dans \emph{son} slot un champ de l'\emph{autre} slot, dans un objet
littéral neuf.
\begin{lstlisting}[language=TypeScript,caption={\file{/tmp/ch15/two_effects.tsx}},label={lst:churn-two-effects}]
import { useState, useEffect } from 'react';
export function TwoEffects() {
  const [a, setA] = useState({ n: 0 });
  const [b, setB] = useState({ n: 0 });
  useEffect(() => { setB({ from: a.n }); }, [a]);
  useEffect(() => { setA({ from: b.n }); }, [b]);
  return <div>{a.n + b.n}</div>;
}
\end{lstlisting}
\end{exemple}

Déroulons ce que fait React. Au montage, les deux effets tournent : \code{setB({...})}
place une référence neuve dans \code{b}, \code{setA({...})} une référence neuve dans
\code{a}. Le composant re-rend. Le premier effet compare son dep \code{a} par
\code{Object.is} avec la valeur précédente : ce sont deux objets distincts, il re-tourne
et écrit un \emph{nouveau} \code{b}. Le second effet fait de même et écrit un nouveau
\code{a}. Rien ne converge jamais : c'est une boucle de rendu infinie, que React
interrompt avec l'erreur \og Maximum update depth exceeded\fg.

\begin{react}[\code{Object.is} sur les deps]
Après un commit, React ré-exécute un effet si l'un de ses deps a changé au sens de
\code{Object.is} par rapport au rendu précédent. Un objet littéral \code{{...}} évalué dans
un effet est une référence neuve à chaque exécution : \code{Object.is} échoue
\emph{toujours}. Une écriture d'état qui stocke une telle référence relance donc tout effet
dont un dep \emph{est} ce slot. Voir \cref{chap:react-modele}.
\end{react}

Voici ce que dit \reactant{} sur ce fichier.

\begin{sortie}
  TwoEffects  (4 hooks)  /tmp/ch15/two_effects.tsx
    error  infinite-loop  [hook:2]  (line 5:2)  these effects form a state-update cycle (`a` → `b` → `a`) where each step stores a fresh reference that re-runs the next effect: infinite render loop
       → a fresh value is written to state `b` here [hook:2] (line 5:20)
       → cycle continues: this effect freshly stores state `a` [hook:3] (line 6:20)
    error  infinite-loop  [hook:3]  (line 6:2)  these effects form a state-update cycle (`a` → `b` → `a`) where each step stores a fresh reference that re-runs the next effect: infinite render loop
       → a fresh value is written to state `a` here [hook:3] (line 6:20)
       → cycle continues: this effect freshly stores state `b` [hook:2] (line 5:20)

⚠  2 error(s) across 1 file(s).
\end{sortie}

Deux \Error, une par effet porteur, et un chemin \code{`a`}~→ \code{`b`}~→ \code{`a`}. C'est le résultat
que ce chapitre explique.

\subsection{Pourquoi les autres bras sont aveugles}

La règle \regle{infinite-loop} (\cref{chap:infinite-loop}) a quatre bras. Deux d'entre
eux ne peuvent \emph{pas} voir ce programme, et comprendre pourquoi motive le graphe.

\paragraph{Le bras point fixe ne voit aucune divergence.} Le point fixe d'un composant
(\cref{chap:point-fixe-composant}) détecte une boucle par le \emph{widening} : un compteur
\code{setN(n + 1)} fait croître l'intervalle jusqu'au seuil, et l'événement de widening
est le signal. Ici la valeur abstraite d'un slot objet est une \emph{stabilité}
(\cref{def:stabilite}) ; la référence stockée est \code{PerRender}, et
\code{PerRender} $\join$ \code{PerRender} $=$ \code{PerRender}. Rien ne croît, rien ne
s'élargit : le point fixe converge en deux tours et ne dit rien. La sémantique de
\emph{valeur} ne distingue pas \og $0 \to 1 \to 0 \to 1$\fg{} de \og atteint
$\{0, 1\}$\fg{} ; le churn est une abstraction de \emph{transition}, et c'est l'argument
d'\adr{42} (contexte) pour ne pas en faire un champ de \code{StateValue}.

\paragraph{Le bras self-churn exige le même slot.} Le bras d'\adr{17} prouve
\code{useEffect(() => setObj({...obj}), [obj])} : le slot écrit \emph{est} un dep du
même effet (must-rerun), l'écriture est sur tous les chemins (must-reach), la valeur est
\code{PerRender} (must-change). Dans l'\cref{ex:churn-two-effects} chaque effet écrit un slot
qu'il ne lit \emph{pas} : \code{setB} sous \code{[a]}, \code{setA} sous \code{[b]}. La
conjonction ne se referme qu'en composant les deux effets.

\subsection{Les trois jambes d'une arête}

Le graphe compose exactement les trois must du bras self-churn, mais \emph{par arête}
plutôt que par effet. Une arête $x \to y$ signifie : \og un changement de $x$ relance un
effet qui stocke une référence fraîche dans $y$\fg{} (doc de module,
\src{src/engine/churn.rs}{3}{6}). Pour qu'elle soit certaine, il faut :
\begin{enumerate}
  \item \textbf{must-rerun} : un dep de l'effet \emph{est} le slot $x$ (ligne
        \relation{effect_triggers} avec \code{exact = true}) ;
  \item \textbf{must-reach} : l'écriture dans $y$ est dans un bloc sur tous les chemins du
        corps (\code{on_all_paths}) ;
  \item \textbf{must-change} : la valeur écrite est \code{Fresh}, c'est-à-dire
        \code{PerRender} (colonne \code{written.fresh} de \relation{slot_writers}).
\end{enumerate}
Un cycle dont \emph{toutes} les arêtes sont certaines est une boucle certaine : c'est la
preuve d'\Error. Une arête où l'une des jambes manque est \code{May} ; un cycle qui en
contient une est un \Warning. Le reste du chapitre déplie cette phrase.

% ---------------------------------------------------------------------------
\section{Le graphe : slots qualifiés, arêtes Must et May}
\label{sec:churn-graphe}

\subsection{Les nœuds : des slots qualifiés}

Un nœud du graphe n'est pas un slot mais un \terme{slot qualifié} : le couple
\code{(ComponentId, HookLabel)}, type \code{QualifiedSlot}
(\srcline{src/ir/types.rs}{9}). La qualification est indispensable dès qu'un effet écrit
le slot d'un \emph{autre} composant à travers un setter passé en prop : la ligne
\relation{slot_writers} porte alors \code{owner = Some(parent)} et son \code{slot} est le
label \emph{du propriétaire} (\cref{def:slot-writer}). La fonction \code{node_of}
(\src{src/engine/churn.rs}{153}{156}) fait cette traduction en une ligne :
\code{(w.owner.unwrap_or(component), w.slot)}.

\subsection{Les types}

\rustsnippet[label={lst:churn-types}]{src/engine/churn.rs}{85}{129}{les quatre types publics du graphe}

Quatre types, dont on retiendra :
\begin{itemize}
  \item \code{EdgeStrength} dérive \code{Ord} avec \code{May < Must} : la déduplication
        garde \og la plus forte\fg{} par simple comparaison ;
  \item \code{ChurnEdge} porte, outre ses extrémités et sa force, le \emph{porteur}
        (\code{component}, \code{effect_label}), le span de l'écriture pour le témoin, et
        deux booléens de \emph{partition} : \code{no_deps} (l'effet n'a pas de tableau de
        deps) et \code{self_slot} (arête guidée par deps d'un slot vers lui-même, réservée au
        bras self-churn) ;
  \item \code{ChurnCycle} est une liste d'indices d'arêtes \emph{en ordre de cycle}
        (\code{edges[i].to == edges[i+1].from}, la dernière rebouclant sur la première) et
        deux faits exacts, \code{all_must} et \code{cross_component} ;
  \item \code{ChurnGraph} est une donnée \emph{programme-entier}, construite une fois par
        programme et partagée par toutes les passes de règles via \code{ProgramRelations}
        (\cref{chap:relations}) ; le compteur \code{BUILDS} (\src{src/engine/churn.rs}{131}{137})
        n'existe qu'en test pour le garantir (\issue{86}).
\end{itemize}

\begin{definition}{Graphe de churn}{churn-graph}
Soit un programme dont les composants sont convergés. Le \emph{graphe de churn} est le
graphe orienté dont les nœuds sont les slots qualifiés $(c, \ell)$ et dont les arêtes
$x \xrightarrow{e} y$ sont portées par un effet $e$ (non mount-only, voir
\cref{sec:churn-sites}) tel que : $e$ possède une ligne d'écriture de $y$ de phase
autre que \code{Handler} et de fraîcheur autre que \code{Not} ; et soit $e$ n'a pas de
tableau de deps et $x = y$ (\emph{auto-arête}), soit une ligne \relation{effect_triggers}
de $e$ nomme $x$. L'arête est \code{Must} quand la ligne de déclencheur est exacte, que
l'écriture est \code{Fresh}, synchrone (\code{block = Some}) et que l'ensemble des blocs
d'écritures fraîches de $y$ dans le corps de $e$ est sur tous les chemins ; elle est
\code{May} sinon. Une arête tuée par la preuve de convergence
(\cref{sec:churn-preuve}) n'existe pas. Un \emph{cycle} est un cycle simple de ce graphe
ne passant par aucune arête \code{self_slot}.
\end{definition}

La définition emprunte deux notions qui seront précisées plus loin (les sites et le kill) ;
pour l'instant, on lit le graphe \emph{sans kill}. Le sens d'une arête est celui de la doc
de module : $x \to y$ dit qu'un \emph{changement} de $x$ relance un effet qui stocke une
référence \emph{fraîche} dans $y$. Le graphe parle donc de transitions, pas de valeurs.

\subsection{La force d'une arête}

\rustsnippet[label={lst:churn-strength}]{src/engine/churn.rs}{530}{543}{la jambe must-reach et la force}

La force se décide en trois conjoints, dans l'ordre de la définition :
\begin{enumerate}
  \item \code{w.written.fresh == Freshness::Fresh} : la colonne \code{written}
        (\cref{chap:relations}) dit que l'argument 0 est une référence \code{PerRender}
        à chaque appel ; \code{Maybe} (valeur opaque, updater imprécis) ne suffit pas ;
  \item \code{w.block.is_some()} : la ligne porte un bloc, ce qui n'arrive que pour une
        écriture \emph{synchrone, une fois par passe} de la région
        (\cref{def:slot-writer}) ; une écriture imbriquée dans une closure, différée,
        de cleanup ou répétée (\code{forEach}) a \code{block = None} et ne peut être que
        \code{May} ;
  \item \code{on_all_paths(f.body_cfg, &fresh_blocks)} : l'ensemble de \emph{tous} les blocs
        d'écritures fraîches du même nœud dans ce corps post-domine l'entrée. Prendre
        l'ensemble plutôt que le seul bloc courant est ce qui rend \code{Must} un
        \code{if (p) setX({...}) else setX({...})} : aucun des deux blocs n'est sur tous les
        chemins, mais leur union l'est.
\end{enumerate}

\rustsnippet[label={lst:churn-on-all-paths}]{src/engine/dominance.rs}{72}{92}{\code{on_all_paths}, la primitive must-reach}

\code{on_all_paths} est un parcours depuis l'entrée qui \emph{évite} les blocs
de l'ensemble : si l'on atteint un bloc sans successeur (une sortie), un chemin s'échappe et
la réponse est \code{false}. Le commentaire parle de BFS, mais la \og file\fg{} est un
\code{Vec} dont on retire le dernier élément (\code{queue.pop()}) : le parcours est en
profondeur. L'ordre de visite est sans effet sur la réponse, qui ne dépend que de
l'ensemble des blocs atteignables. C'est la même primitive que \code{must_on_all_paths} côté
règles (\cref{chap:cfg-analyzer-dominance}), déplacée dans le moteur par \adr{42} \S6
(\og les preuves suivent les données\fg).

La jambe must-rerun, elle, n'apparaît pas dans cet extrait : elle est lue dans les lignes
\relation{effect_triggers} du porteur, partitionnées en \code{exact_local} (le dep
\emph{est} le slot local, \code{exact = true}) et \code{versioned} (le dep est seulement
\emph{versionné} par le slot qualifié, \code{exact = false}) à
\src{src/engine/churn.rs}{313}{323}. Une arête issue de \code{versioned} est \code{May}
quelle que soit sa force d'écriture (\srcline{src/engine/churn.rs}{592}).

\begin{table}[htbp]\centering\small
\begin{tabularx}{\linewidth}{>{\raggedright\arraybackslash}p{42mm}>{\raggedright\arraybackslash}X>{\raggedright\arraybackslash}X}
\toprule
Phase de la ligne & Arête guidée par deps & Auto-arête (effet sans deps) \\
\midrule
\code{Effect} (synchrone, porte un bloc) & \code{Must} si exact $\wedge$ \code{Fresh} $\wedge$ tous les chemins, sinon \code{May} & idem, sans dep à tester \\
\code{Deferred}, \code{Cleanup} & \code{May} & \code{May} (\issue{26}) \\
\code{Handler} & aucune arête & aucune arête \\
\code{Unknown} ($\Top$) & \code{May} & \code{May} \\
\bottomrule
\end{tabularx}
\caption{La table des phases d'\adr{42} \S4 : ce qu'une écriture de corps d'effet peut soutenir selon sa colonne \code{phase}. Le filtre \code{Handler} est à \src{src/engine/churn.rs}{304}{309}, le reste découle de \code{block} (\cref{lst:churn-strength}).}
\label{tab:churn-phases}
\end{table}

Le \cref{tab:churn-phases} n'est pas codé comme une table : il \emph{découle} de deux
faits. Une ligne de phase \code{Handler} est filtrée à la collecte des écritures de l'effet
(\src{src/engine/churn.rs}{304}{309}) : une écriture qu'un listener seul atteint demande un
événement utilisateur par itération, la boucle n'est pas auto-entretenue (preuve par la table
des registrars, \adr{34}). Toute autre ligne entre, et sa colonne \code{block} décide
seule si elle peut être \code{Must}. Le $\Top$ (\code{Unknown}) est donc \code{May} : la
direction \og tirer plus\fg.

\subsection{L'effet sans deps et son auto-arête}

\begin{exemple}{Écriture fraîche sans tableau de deps}{churn-nodeps}
\begin{lstlisting}[language=TypeScript,caption={\file{/tmp/ch15/nodeps.tsx}},label={lst:churn-nodeps}]
import { useState, useEffect } from 'react';
export function NoDeps() {
  const [obj, setObj] = useState({ a: 1 });
  useEffect(() => {
    setObj({ a: 2 });
  });
  return <p>{obj.a}</p>;
}
export function NoDepsThen() {
  const [obj, setObj] = useState({ a: 1 });
  useEffect(() => {
    Promise.resolve().then(() => setObj({ a: 2 }));
  });
  return <p>{obj.a}</p>;
}
\end{lstlisting}
\begin{sortie}
  NoDeps  (2 hooks)  /tmp/ch15/nodeps.tsx
    error  infinite-loop  [hook:1]  (line 4:2)  this effect has no dependency array and stores a fresh reference into state `obj`, so it re-runs after every render and re-triggers itself: infinite render loop
       → a fresh value is written to state `obj` here [hook:1] (line 5:4)
  NoDepsThen  (2 hooks)  /tmp/ch15/nodeps.tsx
    warn   infinite-loop  [hook:1]  (line 11:2)  this effect has no dependency array and may store a fresh reference into state `obj`, so it re-runs after every render and can re-trigger itself: possible infinite render loop
       → a fresh value is written to state `obj` here [hook:1] (line 12:4)
\end{sortie}
\end{exemple}

Un effet sans tableau de deps re-tourne après \emph{chaque} rendu, y compris celui que sa
propre écriture provoque : il n'y a pas de dep à tester, l'auto-arête $y \to y$ est posée
d'office.

\rustsnippet[label={lst:churn-nodeps-edge}]{src/engine/churn.rs}{569}{576}{l'auto-arête d'un effet sans deps}

Dans \code{NoDeps}, l'écriture est synchrone (\code{block = Some(0)}, l'entrée) et
\code{Fresh} : l'arête est \code{Must}, le cycle de longueur 1 est tout-Must, \Error.
Dans \code{NoDepsThen}, la ligne a la phase \code{Deferred} (la table des registrars
connaît \code{.then}) et donc \code{block = None} : l'arête est \code{May}, \Warning.
C'est la cellule \og \code{Deferred} / auto-arête = \code{May}\fg{} du
\cref{tab:churn-phases}, et elle est une décision.

\begin{decision}[ADR-42 \S4]
Avant \adr{42}, le graphe vivait dans \file{rules/helpers} avec sa propre marche des setters,
sans classification de phase : une écriture imbriquée dans \code{.then} était lue \og jamais
auto-entretenue\fg{} et la boucle de \code{NoDepsThen} était silencieuse (\issue{26}, un
FN). Plutôt que de corriger le second walker, \adr{42} le supprime : le graphe devient un
\emph{pli} sur \relation{slot_writers} et \relation{effect_triggers}, et la colonne
\code{phase} décide. Deux cellules de la table changent le comportement : les auto-arêtes
\code{Deferred}/\code{Unknown} apparaissent (côté \og tirer plus\fg) et les arêtes guidées
par deps d'une ligne \code{Handler} disparaissent (appuyées sur la preuve de la table des
registrars). Mesure (\file{docs/precision-log.md}) : 1\,493 $\to$ 1\,494 sur le corpus,
aucune ligne retirée, une ajoutée --- un \Warning{} FP de la cellule $\Top$ (une écriture
sous un \code{useDebouncedCallback} non résolu, dans \code{twenty}) ; la cellule
\issue{26} n'a rien changé sur ce corpus, elle est épinglée par \file{tests/effect_cycles.rs}.
\end{decision}

Le champ \code{no_deps} est \code{deps.list().is_none()}
(\srcline{src/engine/churn.rs}{332}) : un tableau de deps que le moteur ne sait pas lire
(un spread, une variable) est traité comme \emph{absent}, direction \og tirer plus\fg.

\subsection{Deux bras, une relation : \code{self_slot}}

\rustsnippet[label={lst:churn-dep-edges}]{src/engine/churn.rs}{577}{594}{les arêtes guidées par deps, exactes puis versionnées}

Pour chaque slot local exact $x$ des deps, une arête $x \to y$ de la force calculée ; pour
chaque slot qualifié versionné, une arête \code{May} (sauf s'il a déjà été poussé en exact).
Le booléen \code{self_slot} marque l'arête $x \to x$ d'un effet \emph{avec} deps : c'est le
motif \code{setObj({...obj})} sous \code{[obj]} du bras self-churn (\adr{17}), et le graphe
la construit mais ne la cherche jamais dans un cycle.

\begin{decision}[ADR-20 item 2]
Fusionner les deux bras (self-churn et graphe) a été proposé et refusé : le bras self-churn
lit les \emph{membres} du slot (\code{can_retrigger}, bras membre de la preuve) et produit
un diagnostic ancré sur le slot avec sa propre stratification \Error/\Warning/\Info ; le
bras graphe ne le peut pas (savoir si l'écriture de $y$ par $A$ change un dep de $B$ est une
propriété d'une \emph{paire} d'arêtes). La partition est donc explicite : \code{self_slot}
pour \og même slot, avec deps\fg{} (bras self-churn), tout le reste (auto-arêtes sans deps
comprises, poussées avec \code{self_slot = false} dans l'\cref{lst:churn-nodeps-edge}) pour le
graphe. Une relation, deux lecteurs disjoints : aucun double rapport.
\end{decision}

L'arête \code{self_slot} est abandonnée quand \code{write_can_retrigger} répond \code{false}
(\issue{90}, \src{src/engine/churn.rs}{828}{852}) : un updater fonctionnel
\code{prev => ({ ...prev, slug: f(prev) })} laisse chaque membre non nommé
\code{Object.is}-égal, et un dep qui ne lit que ces membres ne peut pas relancer l'effet.
La preuve est syntaxique et \emph{ferme par défaut} : \code{literal_overwrites}
(\src{src/engine/churn.rs}{882}{899}) n'accepte que \code{{ ...prev, a: x, ... }} avec le
spread \emph{en premier} et des clés nommées ; toute autre forme répond \og peut
relancer\fg.

\subsection{Déduplication et déterminisme}

Une clé \code{(from, to, component, effect_label)} identifie une arête. Deux écritures du
même nœud dans le même effet donnent deux candidates ; on garde la plus forte, puis la plus
tôt dans la source (\code{(strength, Reverse(pos))}, \src{src/engine/churn.rs}{545}{566}),
de sorte que la ligne sur laquelle un lecteur s'ancre ne dépende pas de l'ordre de la
relation. Les arêtes sont enfin triées par \code{(component, effect_label, from, to)}
(\src{src/engine/churn.rs}{598}{606}) : \code{result.components} est une \code{HashMap} et
l'ordre de visite varie d'une exécution à l'autre.

\subsection{Trois graphes}

Revenons à l'\cref{ex:churn-two-effects}. La sonde donne les deux lignes d'écriture, les
deux déclencheurs et le graphe (colonnes élaguées) :

\begin{sortie}
-- slot_writers
  slot=0 setter=setA line=Some(6) region=Effect(3) phase=Effect block=Some(0) guard_block=Some(0) fresh=Fresh value.ref=PerRender
  slot=1 setter=setB line=Some(5) region=Effect(2) phase=Effect block=Some(0) guard_block=Some(0) fresh=Fresh value.ref=PerRender
-- effect_triggers
  hook=2 dep=0 slot=(TwoEffects, 0) exact=true
  hook=3 dep=0 slot=(TwoEffects, 1) exact=true
=== churn graph
  edge#0 (TwoEffects, 0) -> (TwoEffects, 1) strength=Must carrier=TwoEffects effect=2 line=Some(5) no_deps=false self_slot=false
  edge#1 (TwoEffects, 1) -> (TwoEffects, 0) strength=Must carrier=TwoEffects effect=3 line=Some(6) no_deps=false self_slot=false
  cycle edges=[0, 1] all_must=true cross_component=false
\end{sortie}

Le dep \code{a} de l'effet 2 \emph{est} le slot 0 (\code{exact=true}) ; l'écriture de
\code{b} est \code{Fresh} et dans le bloc 0, l'entrée, donc sur tous les chemins : arête
$a \to b$ \code{Must}. Symétriquement $b \to a$. La \cref{fig:churn-two-effects} dessine ce
graphe et deux voisins.

\begin{figure}[htbp]\centering
\begin{tikzpicture}[node distance=14mm and 22mm]
  % --- graphe 1 : two_effects
  \begin{scope}[on background layer]
    \node[fill=ra-bg,rounded corners=6pt,fit={(-1,-0.9) (2.6,0.9)},inner sep=0pt] {};
  \end{scope}
  \node[noeud] (a1) at (0,0) {$a$};
  \node[noeud] (b1) at (1.8,0) {$b$};
  \draw[flechemust] (a1) to[bend left=30] node[etiquette,above] {\code{e2}} (b1);
  \draw[flechemust] (b1) to[bend left=30] node[etiquette,below] {\code{e3}} (a1);
  \node[etiquette] at ($(a1)!0.5!(b1) + (0,-1.3)$) {\code{TwoEffects} : une SCC tout-Must, \Error};
  % --- graphe 2 : three
  \begin{scope}[xshift=52mm]
  \begin{scope}[on background layer]
    \node[fill=ra-bg,rounded corners=6pt,fit={(-1,-0.9) (4.4,1.9)},inner sep=0pt] {};
  \end{scope}
  \node[noeud] (a2) at (0,0) {$a$};
  \node[noeud] (b2) at (1.8,0) {$b$};
  \node[noeud] (c2) at (3.6,0) {$c$};
  \draw[flechemust] (a2) -- node[etiquette,above] {\code{e3}} (b2);
  \draw[flechemust] (b2) -- node[etiquette,above] {\code{e4}} (c2);
  \draw[flechemay] (c2) to[bend right=45] node[etiquette,above] {\code{e5}, sous \code{c.n > 3}} (a2);
  \node[etiquette,below=9mm of b2] {\code{Three} : une SCC, une arête May, \Warning};
  \end{scope}
  % --- graphe 3 : dag
  \begin{scope}[yshift=-32mm]
  \node[noeud] (a3) at (0,0) {$a$};
  \node[noeud] (b3) at (1.8,0) {$b$};
  \node[noeud] (c3) at (3.6,0) {$c$};
  \draw[flechemust] (a3) -- node[etiquette,above] {\code{e3}} (b3);
  \draw[flechemust] (b3) -- node[etiquette,above] {\code{e4}} (c3);
  \node[etiquette,below=7mm of b3] {\code{Dag} : trois SCC triviales, aucun cycle ; \code{onClick} ne porte rien};
  \end{scope}
  % légende
  \begin{scope}[xshift=62mm,yshift=-32mm]
  \draw[flechemust] (0,0.3) -- (1,0.3) node[etiquette,right,text=black] {arête \code{Must}};
  \draw[flechemay] (0,-0.3) -- (1,-0.3) node[etiquette,right,text=black] {arête \code{May}};
  \node[fill=ra-bg,rounded corners=3pt,minimum width=9mm,minimum height=5mm] at (0.5,-0.95) {};
  \node[etiquette,right,text=black] at (1,-0.95) {SCC cyclique};
  \end{scope}
\end{tikzpicture}
\caption{Trois graphes de churn (\file{two_effects.tsx}, \file{three.tsx}, \file{dag.tsx}). Les nœuds sont des slots qualifiés du même composant ; \code{eN} est le label du hook porteur. Une arête pleine est \code{Must}, une arête pointillée \code{May}.}
\label{fig:churn-two-effects}
\end{figure}

\begin{exemple}{Un cycle à trois, une garde}{churn-three}
\begin{lstlisting}[language=TypeScript,caption={\file{/tmp/ch15/three.tsx}},label={lst:churn-three}]
import { useState, useEffect } from 'react';
export function Three() {
  const [a, setA] = useState({ n: 0 });
  const [b, setB] = useState({ n: 0 });
  const [c, setC] = useState({ n: 0 });
  useEffect(() => { setB({ from: a.n }); }, [a]);
  useEffect(() => { setC({ from: b.n }); }, [b]);
  useEffect(() => { if (c.n > 3) setA({ from: c.n }); }, [c]);
  return <div>{a.n + b.n + c.n}</div>;
}
\end{lstlisting}
La sonde montre \code{block=Some(1)} pour \code{setA} (le bloc de la branche vraie, pas
l'entrée) : \code{on_all_paths} échoue, l'arête $c \to a$ est \code{May}, le cycle n'est
pas \code{all_must}. Trois \Warning{} \og may form a state-update cycle
(\code{`a`}~→ \code{`b`}~→ \code{`c`}~→ \code{`a`})\fg, un par effet porteur. Le même programme sans la garde
(\file{three_must.tsx}) donne trois \Error. Le lecteur notera que la garde \code{c.n > 3}
ne \emph{meurt} pas sous l'écriture \code{setA({...})} : elle teste \code{c}, pas
\code{a} ; le kill de convergence de la \cref{sec:churn-preuve} n'a rien à dire ici.
\end{exemple}

\begin{exemple}{Une chaîne dérivée est un DAG}{churn-dag}
Dans \file{dag.tsx}, $a \to b \to c$ sans retour, et \code{setA} n'est appelé que dans
\code{onClick}. Le graphe a deux arêtes \code{Must} et aucun cycle ; la ligne de
\code{setA} a \code{region=Handler(5) phase=Handler} et n'entre pas. La règle n'émet que
deux \Info{} (bras self-churn : \og freshly recreates state \code{`b`} outside its deps.
No update cycle was found\fg), le marqueur d'imprécision résiduelle de la
\cref{sec:churn-certification}.
\end{exemple}

% ---------------------------------------------------------------------------
\section{Trouver les cycles : Tarjan en deux passes}
\label{sec:churn-tarjan}

\subsection{Pourquoi deux passes}

Le graphe une fois construit, il faut en extraire les cycles et leur attribuer une force.
Une idée naïve serait de chercher les cycles dans le graphe entier et de déclarer
\code{all_must} ceux dont toutes les arêtes sont \code{Must}. Elle rate un cas : un cycle
tout-Must $a \to b \to a$ noyé dans une composante fortement connexe plus grande qui
contient aussi une arête \code{May} $b \to c \to a$. Un seul cycle simple étant reconstruit
par composante, on risquerait de ne rapporter que le \code{May}. D'où deux passes.

\rustsnippet[label={lst:churn-find-cycles}]{src/engine/churn.rs}{635}{662}{\code{find_cycles} : sous-graphe Must d'abord, puis tout}

\begin{enumerate}
  \item \textbf{Passe Must.} On ne garde que les arêtes \code{Must} non \code{self_slot}
        (\code{must_idx}) et l'on cherche les composantes fortement connexes (SCC)
        cycliques ; chacune donne un cycle \code{all_must = true}, et ses nœuds sont
        marqués \code{covered_nodes}.
  \item \textbf{Passe complète.} On reprend toutes les arêtes non \code{self_slot}
        (\code{all_idx}) ; un cycle dont un nœud est déjà couvert est \emph{sauté} : l'\Error{}
        signale déjà cette région, on ne rapporte pas un cycle plus faible qui la recoupe. Les
        autres donnent des cycles \code{all_must = false}.
\end{enumerate}

\begin{piege}
La seconde passe saute un cycle \code{May} dès qu'\emph{un} de ses nœuds est dans une SCC
tout-Must. Un cycle $a \to b \to a$ tout-Must et un cycle $b \to d \to b$ \code{May}
partagent $b$ : seul le premier est rapporté. Ce n'est pas un FN, la boucle certaine
couvre la région ; mais un lecteur qui compte les cycles de \code{graph.cycles} doit savoir
qu'il ne compte pas tous les cycles du graphe.
\end{piege}

\subsection{Composantes et reconstruction}

\code{cycles_in} (\src{src/engine/churn.rs}{680}{713}) construit une table de nœuds
\emph{triée} (déterminisme), une liste d'adjacence nœud $\to$ indices d'arêtes, puis
appelle \code{tarjan_sccs} (\src{src/engine/churn.rs}{716}{779}), un Tarjan récursif
classique — \og les graphes sont minuscules\fg, dit le commentaire. Une SCC est cyclique
si elle a au moins deux nœuds ou si un de ses nœuds porte une auto-arête
(\src{src/engine/churn.rs}{700}{703}). Pour chaque SCC cyclique, \code{find_cycle_from}
(\src{src/engine/churn.rs}{783}{814}) reconstruit \emph{un} cycle simple par un parcours en
profondeur itératif depuis le plus petit nœud de la composante, en restant dans la SCC : il
existe toujours, et son chemin d'arêtes devient \code{edge_idx}. Enfin \code{make_cycle}
(\src{src/engine/churn.rs}{664}{676}) pose \code{cross_component} : l'ensemble des
composants impliqués (propriétaires de \code{from}, de \code{to}, et porteurs) a plus d'un
élément.

Le coût est linéaire : $O(V + E)$ pour Tarjan, un parcours par SCC pour la reconstruction,
le tout \emph{une fois par programme}.

\subsection{Ce que prouve un cycle tout-Must}

\begin{lemme}{Composition des arêtes Must}{churn-must-cycle}
Soit $x_0 \to x_1 \to \dots \to x_{n-1} \to x_0$ un cycle du graphe dont toutes les arêtes
sont \code{Must}, portées par des effets $e_0, \dots, e_{n-1}$ d'un même composant. Dès que
le composant est monté et que $e_0$ s'exécute une fois, chaque tour de rendu ré-exécute
$e_0, \dots, e_{n-1}$ et pose une référence neuve dans chacun de $x_1, \dots, x_0$ : la
boucle de rendu est infinie.
\end{lemme}
\begin{proof}
Par récurrence sur les tours. Au montage, React exécute tous les effets, donc $e_0$
s'exécute. Supposons $e_i$ exécuté au tour $t$. Son écriture de $x_{i+1}$ est sur tous les
chemins du corps (must-reach) : elle a lieu. Sa valeur est \code{PerRender} (must-change) :
c'est une référence que \code{Object.is} distingue de la précédente, React planifie un
rendu. Au tour $t+1$, un dep de $e_{i+1}$ \emph{est} le slot $x_{i+1}$ (must-rerun) :
$e_{i+1}$ s'exécute. La chaîne se referme sur $e_0$ au tour $t+n$.
\end{proof}

C'est la version composée du triple must d'\adr{17} \S3, et l'argument de soundness 1
d'\adr{18}. Le lemme énonce une propriété du programme \emph{concret} : c'est pourquoi la
règle peut frapper une \Error{} (\og défaut certain dès que le code s'exécute\fg). Deux
réserves, traitées plus loin : la preuve suppose que les jambes sont vraiment must (la
\cref{sec:churn-cross} montre pourquoi c'est faux au travers d'un prop), et qu'aucune
garde ne meurt (\cref{sec:churn-preuve}).

Voici la sortie observée pour \file{three_must.tsx}, la version sans garde de
l'\cref{ex:churn-three} :

\begin{sortie}
  ThreeMust  (6 hooks)  three_must.tsx
    error  infinite-loop  [hook:3]  (line 6:2)  these effects form a state-update cycle (`a` → `b` → `c` → `a`) where each step stores a fresh reference that re-runs the next effect: infinite render loop
       → a fresh value is written to state `b` here [hook:3] (line 6:20)
       → cycle continues: this effect freshly stores state `c` [hook:4] (line 7:20)
       → cycle continues: this effect freshly stores state `a` [hook:5] (line 8:20)
    error  infinite-loop  [hook:4]  (line 7:2)  these effects form a state-update cycle (`a` → `b` → `c` → `a`) …
    error  infinite-loop  [hook:5]  (line 8:2)  these effects form a state-update cycle (`a` → `b` → `c` → `a`) …

⚠  3 error(s) across 1 file(s).
\end{sortie}

Le chemin affiché part toujours du plus petit nœud de la SCC (\code{find_cycle_from} démarre
en \code{scc.iter().min()}), quel que soit l'effet porteur du diagnostic : les trois
messages disent \code{`a`}~→ \code{`b`}~→ \code{`c`}~→ \code{`a`}, et seule la première note (\og a fresh value
is written … here\fg) change.

% ---------------------------------------------------------------------------
\section{Le kill de convergence, première époque : un seul écrivain}
\label{sec:churn-kill-histoire}

\subsection{Le motif \og fetch once\fg{} en paire}

Le graphe de la section précédente a un défaut immédiat : il voit des boucles là où le
programme converge. Le motif le plus courant est la garde \og si vide, remplir\fg.

\begin{exemple}{Deux chargements gardés}{churn-fetch-once}
\begin{lstlisting}[language=TypeScript,caption={\file{/tmp/ch15/fetch_once_pair.tsx}},label={lst:churn-fetch-once}]
import { useState, useEffect } from 'react';
export function FetchOncePair() {
  const [a, setA] = useState(null);
  const [b, setB] = useState(null);
  useEffect(() => { if (!b) setB({ src: 'e1' }); }, [a]);
  useEffect(() => { if (!a) setA({ src: 'e2' }); }, [b]);
  return <div>{a && b ? 'ok' : 'loading'}</div>;
}
\end{lstlisting}
\begin{sortie}
  FetchOncePair  (4 hooks)  /tmp/ch15/fetch_once_pair.tsx  ✓
    verified  infinite-loop  no effect diverges into an infinite render loop
\end{sortie}
\end{exemple}

Concrètement : au montage \code{b} est \code{null}, le premier effet écrit \code{{...}}
dans \code{b} ; \code{a} est \code{null}, le second écrit \code{{...}} dans \code{a}. Au
rendu suivant, les deux effets re-tournent (leurs deps ont changé) mais les deux gardes sont
fausses : plus aucune écriture, le composant est stable en deux rendus. Sans kill, le graphe
aurait pourtant deux arêtes \code{May} (les écritures sont sous une garde,
\code{block=Some(1)} n'est pas sur tous les chemins) et un cycle $a \to b \to a$ : un
\Warning{} faux positif sur l'idiome le plus banal de React. La sonde confirme que le graphe
est \emph{vide} :

\begin{sortie}
-- slot_writers
  slot=0 setter=setA line=Some(6) region=Effect(3) phase=Effect block=Some(1) guard_block=Some(1) fresh=Fresh value.ref=PerRender
  slot=1 setter=setB line=Some(5) region=Effect(2) phase=Effect block=Some(1) guard_block=Some(1) fresh=Fresh value.ref=PerRender
-- effect_triggers
  hook=2 dep=0 slot=(FetchOncePair, 0) exact=true
  hook=3 dep=0 slot=(FetchOncePair, 1) exact=true
=== churn graph
\end{sortie}

\subsection{Pourquoi la garde meurt}

L'argument est celui de la preuve mono-site du \cref{chap:relations}
(\code{converges_once_written}, section sur la preuve mono-site), que nous rappelons
sur ce cas. Le corps du premier effet est abaissé en un bloc 0 terminé par
\code{Branch { cond: !b }}, dont l'arête \code{IfTrue} mène au bloc 1 qui contient
\code{setB({...})}. \code{guard_chain(cfg, 1)} remonte la chaîne des prédécesseurs uniques :
$[1, 0]$. \code{site_guards} lit la condition et la polarité de l'arête : $(\code{!b},
\code{true})$, que \code{expand_guard} réécrit en $(\code{b}, \code{false})$. Le bras
\emph{valeur} de \code{dead_once_written} ré-lie alors \code{b} à la valeur écrite --- une
référence \code{PerRender}, donc \emph{truthy} et non nulle --- et applique le
rétrécissement de branche \og \code{b} falsy\fg{} : $\Bot$. La garde est morte sous sa
propre écriture : une fois \code{{...}} posé dans \code{b}, ce site ne peut plus tirer.

\begin{remarque}
Le kill porte sur le slot \emph{écrit}, pas sur le slot lu par les deps. La garde
\code{!b} du premier effet parle de \code{b}, que ce même effet écrit ; le dep \code{a}
est sans importance pour la preuve. C'est ce qui permet de tuer l'arête $a \to b$ : elle
est \og une écriture de $b$ qui ne tire qu'une fois\fg, d'où qu'elle soit déclenchée.
\end{remarque}

\subsection{Le reviver : une écriture voisine ranime la garde}

\begin{exemple}{Un troisième effet remet \code{null}}{churn-multi-writer}
\begin{lstlisting}[language=TypeScript,caption={\file{/tmp/ch15/multi_writer.tsx}},label={lst:churn-multi-writer}]
import { useState, useEffect } from 'react';
export function MultiWriter() {
  const [a, setA] = useState(null);
  const [b, setB] = useState(null);
  useEffect(() => { if (!b) setB({ src: 'e1' }); }, [a]);
  useEffect(() => { setA({ src: 'e2' }); }, [b]);
  useEffect(() => { setB(null); }, [a]);
  return <div/>;
}
\end{lstlisting}
\begin{sortie}
  MultiWriter  (5 hooks)  /tmp/ch15/multi_writer.tsx
    warn   infinite-loop  [hook:2]  (line 5:2)  these effects may form a state-update cycle (`a` → `b` → `a`) where each step may store a fresh reference that re-runs the next effect: possible infinite render loop
       → a fresh value is written to state `b` here [hook:2] (line 5:28)
       → cycle continues: this effect freshly stores state `a` [hook:3] (line 6:20)
    warn   infinite-loop  [hook:3]  (line 6:2)  these effects may form a state-update cycle (`a` → `b` → `a`) …
\end{sortie}
\end{exemple}

La garde \code{!b} meurt toujours sous sa propre écriture. Mais le troisième effet,
lui aussi relancé par \code{a}, remet \code{null} dans \code{b} au tour automatique
suivant : la garde \emph{tient de nouveau}, le premier effet peut réécrire \code{{...}}, et
la boucle vit. La sonde le montre : trois lignes d'écriture, dont \code{setB(null)} avec
\code{fresh=Not value.null=true}, et le graphe garde ses deux arêtes.

\begin{sortie}
  slot=0 setter=setA line=Some(6) region=Effect(3) phase=Effect block=Some(0) fresh=Fresh value.ref=PerRender
  slot=1 setter=setB line=Some(5) region=Effect(2) phase=Effect block=Some(1) fresh=Fresh value.ref=PerRender
  slot=1 setter=setB line=Some(7) region=Effect(4) phase=Effect block=Some(0) fresh=Not value.ref=Bottom value.null=true
=== churn graph
  edge#0 (MultiWriter, 0) -> (MultiWriter, 1) strength=May carrier=MultiWriter effect=2 line=Some(5) no_deps=false self_slot=false
  edge#1 (MultiWriter, 1) -> (MultiWriter, 0) strength=Must carrier=MultiWriter effect=3 line=Some(6) no_deps=false self_slot=false
  cycle edges=[0, 1] all_must=false cross_component=false
\end{sortie}

L'écriture \code{setB(null)} n'est pas fraîche : elle ne porte \emph{aucune} arête (le
\code{continue} de \src{src/engine/churn.rs}{503}{504}). Elle compte pourtant : c'est un
\terme{reviver}, un site qui ranime une garde qu'une autre écriture avait éteinte. Une
preuve de convergence qui l'ignorerait serait insound.

\subsection{La réponse d'ADR-018 : un seul écrivain}

\begin{decision}[ADR-18]
Le graphe de churn naît le 16 juillet avec un kill de convergence \emph{conditionnel} : la
preuve mono-site est appliquée à une arête seulement si le slot écrit a \textbf{un unique
site d'écriture d'effet programme-entier}, quelle que soit la fraîcheur des autres (une
écriture stable ranime aussi une garde). Les handlers ne sont pas comptés (ils demandent un
événement). Argument de soundness (\adr{18}, point 2) : le kill est appliqué strictement
\emph{moins} souvent qu'une preuve naïve, et retirer un kill ne peut qu'ajouter des
arêtes, jamais perdre un cycle réel. Prix : un slot à plusieurs écrivains garde toutes ses
arêtes, au pire un \Warning{} faux positif (\issue{39}).
\end{decision}

Le faux positif de \issue{39} est réel et fréquent : deux effets qui chargent chacun
\emph{le même} slot derrière la \emph{même} garde.

\begin{exemple}{Deux écrivains gardés qui s'éteignent l'un l'autre}{churn-settle}
\begin{lstlisting}[language=TypeScript,caption={\file{/tmp/ch15/settle_each_other.tsx}},label={lst:churn-settle}]
import { useState, useEffect } from 'react';
export function SettleEachOther() {
  const [s, setS] = useState(null);
  const [b, setB] = useState(null);
  useEffect(() => { if (!s) setS({ from: 'e1' }); }, [b]);
  useEffect(() => { if (!s) setS({ from: 'e2' }); }, [b]);
  useEffect(() => { setB({ from: s }); }, [s]);
  return <div>{s ? 'y' : 'n'}</div>;
}
\end{lstlisting}
Quelle que soit l'écriture qui tire la première, \code{s} devient truthy et \emph{les deux}
gardes meurent. Sous \adr{18}, le slot \code{s} a deux écrivains : aucun kill, deux
arêtes \code{May} $b \to s$, une arête \code{Must} $s \to b$, un cycle, deux \Warning.
Au commit \snapshotCommit, la sonde ne garde que l'arête $s \to b$ (aucun cycle) et la
règle n'émet qu'un \Info{} sur l'écriture de \code{b} hors deps. C'est le programme du test
\code{two_guarded_writers_of_one_slot_}\allowbreak\code{that_settle_each_others_guard_converge}
de \file{tests/effect_cycles.rs}.
\end{exemple}

Résoudre \issue{39} exige de poser la \emph{bonne} question : non pas \og combien
d'écrivains ?\fg{} mais \og sous \emph{chaque} écriture du slot, la garde meurt-elle ?\fg.
C'est la preuve multi-site de la \cref{sec:churn-preuve}. Avant de l'écrire, il faut
savoir de quelles écritures on parle.

% ---------------------------------------------------------------------------
\section{Qu'est-ce qu'un site ?}
\label{sec:churn-sites}

\subsection{Qui peut ranimer une garde ?}

\adr{18} ne comptait que les \emph{écritures d'effet}. Mais tout ce que la boucle
automatique exécute peut ranimer une garde : une écriture pendant le rendu, une écriture
dans un \code{useMemo}, et même l'effet de montage d'un \emph{enfant} que la boucle démonte
et remonte. Trois faux négatifs de \issue{162}.

\begin{exemple}{Une écriture au rendu ranime une garde d'effet}{churn-render-write}
\begin{lstlisting}[language=TypeScript,caption={\file{/tmp/ch15/render_write.tsx}},label={lst:churn-render-write}]
import { useState, useEffect } from 'react';
export function RenderWrite({ cond }: { cond: boolean }) {
  const [s, setS] = useState(null);
  if (cond && s) setS(null);
  useEffect(() => { if (!s) setS({ fresh: true }); }, [s]);
  return <div>{s ? 'y' : 'n'}</div>;
}
\end{lstlisting}
Avec \code{cond} vrai : \code{{...}} $\to$ rendu $\to$ \code{null} $\to$ effet $\to$
\code{{...}} $\to$ … La sonde montre la ligne \code{region=Render phase=Render
block=Some(3) fresh=Not value.null=true} à côté de la ligne d'effet, et l'arête
\code{self_slot} $s \to s$ \code{May} survit ; la règle émet un \Warning{}
\og may store a fresh reference into state \code{`s`} which its deps react to\fg. Avant
\snapshotCommit, la liste des sites n'était construite que sur les lignes d'effet : l'écriture
du rendu ne ranimait rien et l'arête était tuée. Test :
\code{a_render_phase_write_revives_an_effect_guard}.
\end{exemple}

\begin{definition}{Site de convergence}{churn-site}
Un \emph{site} (au sens du graphe de churn) est une ligne \relation{slot_writers} d'un
composant, de phase autre que \code{Handler}, prise avec le corps (CFG) dans lequel elle
s'exécute :
\begin{itemize}
  \item région \code{Render} : le CFG de rendu ;
  \item région \code{Effect(l)} d'un effet \emph{non} mount-only, ou \code{Memo(l)} : le
        corps du hook $l$ ;
  \item région \code{Effect(l)} d'un effet \emph{mount-only} (deps d'arité exactement 0) :
        seulement si la ligne est \emph{étrangère} (\code{owner = Some}) \emph{et} que le
        propriétaire ne garde pas ce composant monté à travers la boucle
        (\code{stays_mounted} faux) ; le corps du hook $l$ ;
  \item régions \code{Callback(_)} et \code{Handler(_)} : jamais (les écritures d'un
        callback sont des lignes du corps qui l'appelle, la marche les a réifiées).
\end{itemize}
L'ensemble \code{foreign} est celui des slots qualifiés qu'un site d'un \emph{autre}
composant écrit.
\end{definition}

Ce \og site\fg{} n'est pas le site d'écriture de l'IR (\cref{def:site}, un \code{ExprId}) :
c'est une ligne de relation \emph{située} dans un corps, l'unité sur laquelle la preuve de
convergence raisonne. Le code qui les collecte suit la définition mot pour mot.

\rustsnippet[label={lst:churn-sites}]{src/engine/churn.rs}{259}{287}{la collecte des sites}

Deux filtres distincts jouent ici, et il faut les distinguer (\cref{chap:relations},
région $\neq$ phase) : le filtre sur la \emph{phase} (\code{!= Handler}) écarte une écriture
qu'un listener non réifié atteint, où qu'elle soit écrite ; le filtre sur la \emph{région}
écarte les lignes de région \code{Callback}/\code{Handler} (un \code{setTimeout} dans un
\code{onClick} est de phase \code{Deferred} mais de région \code{Handler} : exclu par la
région). Le \code{expect("every effect write is a site")} de
\srcline{src/engine/churn.rs}{517} est l'invariant qui en découle : toute écriture non-handler
d'un effet non mount-only est dans \code{all_sites}.

\subsection{L'effet de montage d'un enfant}

Un effet à deps \code{[]} tourne une fois par \emph{montage}. Pour le composant lui-même,
qui reste monté pendant la boucle, il ne peut rien ranimer indéfiniment : ses lignes ne sont
pas des sites (cas \code{None => None}). Mais un enfant que le parent rend
\emph{conditionnellement} est démonté et remonté à chaque tour de la boucle, et son effet
\code{[]} tire à chaque tour.

\begin{exemple}{Un enfant qui remonte}{churn-remount}
\begin{lstlisting}[language=TypeScript,caption={\file{/tmp/ch15/remount.tsx}},label={lst:churn-remount}]
import { useState, useEffect } from 'react';
function Child({ onReady }: { onReady: (v: any) => void }) {
  useEffect(() => { onReady(null); }, []);
  return <span />;
}
export function Parent() {
  const [s, setS] = useState(null);
  useEffect(() => { if (!s) setS({ fresh: true }); }, [s]);
  return <div>{s ? <Child onReady={setS} /> : null}</div>;
}
\end{lstlisting}
\begin{sortie}
  Child  (1 hooks)  /tmp/ch15/remount.tsx  ✓
  Parent  (2 hooks)  /tmp/ch15/remount.tsx
    warn   infinite-loop  [hook:0]  (line 8:2)  this effect may store a fresh reference into state `s` which its deps react to: possible infinite render loop
       → a fresh value is written to state `s` here [hook:1] (line 8:28)
\end{sortie}
\code{s} vide $\to$ le parent écrit \code{{...}} $\to$ \code{<Child/>} monte $\to$ son
effet écrit \code{null} dans le slot du parent $\to$ \code{<Child/>} démonte, \code{s} vide,
et ainsi de suite. La ligne de \code{Child} est étrangère (\code{owner=Some("Parent")},
\code{fresh=Not}) ; \code{stays_mounted(Parent)} échoue car l'élément \code{<Child/>}
est sous la garde \code{s} --- le ternaire JSX est abaissé en \code{Branch} et l'élément
vit dans le bloc \code{IfTrue} (\code{lower_ternary},
\src{src/lowering/expr_lower.rs}{675}{724}) --- et \code{s}, un état, ne \emph{tient} pas ; la ligne est un site, donc
\code{(Parent, 0)} est dans \code{foreign} et la preuve pour \code{setS({...})} répond
\code{false} : l'arête \code{self_slot} reste, le bras self-churn émet un \Warning. Test :
\code{a_child_the_loop_remounts_}\allowbreak\code{fires_its_mount_effect_every_round}.
\end{exemple}

La fonction \code{stays_mounted} (\src{src/engine/churn.rs}{374}{430}) répond
\og oui\fg{} seulement si le parent rend au moins un élément de l'enfant, et si
\emph{chaque} occurrence est hors de toute closure (un \code{items.map(i => <Child/>)}
monte un enfant par élément que la boucle peut ajouter), sous des gardes qui \emph{tiennent}
(\code{Invariance::holds}, \cref{sec:churn-preuve}) et sans \code{key}, ou avec une
\code{key} qui tient. Un parent qui ne rend aucun tel élément --- l'enfant est atteint via
un composant intermédiaire --- \emph{échoue fermé}. Deux variantes du test le montrent :
\file{stays_mounted.tsx}, où \code{<Child/>} est derrière \code{if (!ready) return null}
avec \code{ready} un prop, est silencieuse (\og verified\fg) ; \file{keyed.tsx}, où
\code{<Child key={s ? 'a' : 'b'}/>} est rendu sur tous les chemins mais avec une clé qui
lit le slot, est un \Warning{} : React remonte l'enfant à chaque écriture.

\begin{react}[Montage, \code{key}, effets \code{[]}]
React identifie un enfant par sa position dans l'arbre et sa \code{key}. Si l'élément
disparaît d'un rendu (branche \code{null}) ou change de \code{key}, l'instance est démontée
(ses cleanups tournent) et une nouvelle est montée au rendu suivant : tous ses effets,
\code{[]} compris, s'exécutent de nouveau. Un effet \og mount-only\fg{} ne l'est donc que
pour une instance qui reste montée.
\end{react}

\begin{piege}
L'étiquette du diagnostic ci-dessus est \code{[hook:0]} alors que l'effet a le label 1 :
le bras self-churn ancre son diagnostic sur le \emph{slot} (\code{.with_label(state_label)},
\srcline{src/rules/impls/infinite_loop.rs}{524}) et seule l'étape \code{Write} du témoin
porte le label de l'effet. Le bras graphe, lui, ancre sur l'effet porteur.
\end{piege}

\subsection{La vue du site pour la preuve : \code{WriteSite}}

La preuve ne lit pas les lignes brutes mais une projection, \code{WriteSite}, construite par
\code{site_of} (\src{src/engine/churn.rs}{431}{443}).

\rustsnippet[label={lst:churn-writesite}]{src/engine/guards.rs}{92}{121}{\code{WriteSite}, la vue d'un site pour la preuve}

Champ par champ :
\begin{description}
  \item[\code{component}, \code{cfg}, \code{slot}] le composant dont le corps contient
        l'écriture, ce corps, et le slot \emph{qualifié} écrit (\code{node_of}).
  \item[\code{guard_block}] le bloc dont les gardes dominantes conditionnent l'écriture :
        son propre bloc pour une écriture synchrone, \emph{l'instruction qui l'a planifiée}
        pour une écriture imbriquée, différée ou répétée (une continuation \code{.then} ne
        tourne que si la passe qui l'a planifiée a pris toutes les branches au-dessus,
        \issue{160}). C'est la colonne \code{guard_block} de la ligne
        (\cref{def:slot-writer}).
  \item[\code{block}] le bloc de l'écriture quand elle est synchrone, une fois par passe : ce
        qui la fait \emph{co-exécuter} avec les sites qu'elle domine. \code{None} sinon.
  \item[\code{bounded}] \code{phase != Unknown} : l'écriture tourne au plus une fois par
        exécution de l'instruction qui la planifie, ou sur son propre timer (les ticks
        suivants d'un \code{setInterval} sont des événements externes, pas des tours de
        boucle). Un callback confié à un callee inconnu peut être rappelé indéfiniment : ses
        gardes disent quand il a été \emph{remis}, pas combien de fois il tourne.
  \item[\code{convergent}] déjà prouvé \og tire au plus une fois dans la boucle
        automatique\fg{} : il ne ranime rien indéfiniment (\cref{sec:churn-point-fixe}).
  \item[\code{value}, \code{expr}] la valeur abstraite laissée dans le slot et l'argument 0
        tel qu'écrit (colonne \code{written}).
\end{description}

L'ensemble \code{foreign} (\src{src/engine/churn.rs}{344}{350}) se déduit des sites : un
slot qualifié dont le propriétaire n'est pas le composant du site. Un tel slot n'est
\emph{jamais} tué : les gardes de l'autre site vivent dans des corps que l'environnement de
ce composant ne sait pas lire.

% ---------------------------------------------------------------------------
\section[La preuve multi-site]{La preuve multi-site : \code{converges_under_all_writes}}
\label{sec:churn-preuve}

\subsection{La question, en une phrase}

Un site \emph{converge} s'il tire au plus une fois dans la boucle automatique --- la suite
de rendus que les seules écritures d'état entretiennent, sans événement utilisateur. Pour
le prouver, ses gardes doivent mourir :
\begin{enumerate}
  \item[(a)] sous \emph{l'ensemble tueur} : sa propre écriture et toute écriture synchrone
        d'un slot local située sur sa chaîne de gardes, car ces écritures s'exécutent à
        \emph{chaque} passe où lui-même s'exécute ;
  \item[(b)] sous \emph{chaque autre écriture vivante} des slots de cet ensemble, prise avec
        les \emph{faits invariants} sous lesquels cet autre site a tiré.
\end{enumerate}
La condition (a) généralise la preuve mono-site ; la condition (b) est ce qui manquait à
\adr{18} : au lieu de compter les écrivains, on les \emph{interroge} un par un.

\begin{definition}{Ensemble tueur, reviver vivant, faits invariants}{churn-kill-set}
Soit $p$ un site de composant $C$, de bloc de garde $g$ et de chaîne
$\mathrm{chain}(g)$ (les blocs à prédécesseur unique au-dessus de $g$, lui compris).
\begin{itemize}
  \item L'\emph{ensemble tueur} de $p$ est la famille de réécritures
        $\{\ell \mapsto v\}$ formée de la propre écriture de $p$ et de celle de tout site $q$
        du \emph{même corps}, écrivant un slot \emph{local} de $C$, dont \code{block} est dans
        $\mathrm{chain}(g)$. Deux écritures du même slot donnent le join de leurs valeurs.
  \item Un \emph{reviver vivant} d'un slot $\ell$ de l'ensemble tueur est un site $q \neq p$
        de $C$, hors ensemble tueur, écrivant $(C, \ell)$, et non déjà prouvé convergent.
  \item Les \emph{faits} d'un site $q$ sont les conjoints de ses gardes qui
        \emph{tiennent} à travers la boucle (\cref{def:churn-invariance}).
\end{itemize}
$p$ \emph{converge} si ses gardes sont mortes sous l'ensemble tueur et, pour chaque reviver
vivant $q$ d'un slot $\ell$ de l'ensemble tueur, mortes sous l'environnement rétréci par
les faits de $q$ dans lequel $\ell$ reçoit la valeur écrite par $q$ et les autres slots de
l'ensemble tueur gardent leur valeur tueuse s'ils n'ont aucun reviver vivant, $\Top$ sinon.
Si un slot de l'ensemble tueur est dans \code{foreign}, $p$ ne converge pas.
\end{definition}

\subsection{Première moitié : préconditions et ensemble tueur}

\rustsnippet[label={lst:churn-proof-1}]{src/engine/guards.rs}{137}{197}{\code{converges_under_all_writes} : signature et première moitié}

La signature dit déjà tout le contrat : le site est désigné par son indice \code{i} dans
\code{peers} (les sites d'\emph{un} composant), sa valeur propre \code{own_value} est
fournie par l'appelant, et \code{foreign}, \code{invariance}, \code{props_hold} sont les
faits de contexte calculés une fois par \code{build_edges}. Les préconditions \emph{échouent fermées} : pas de \code{guard_block} (ligne sans position
plaçable), site non borné (\code{Unknown} : un callee inconnu peut rappeler le callback),
slot étranger (\code{site.slot.0 != site.component}), aucune garde --- dans les quatre cas la
réponse est \code{false} et l'arête reste. Puis :
\begin{itemize}
  \item \code{co} : les pairs du même corps (\code{ptr::eq(p.cfg, site.cfg)}), d'un slot
        local, dont \code{block} est sur la chaîne. Ils \emph{co-exécutent} avec le site :
        toute passe qui atteint $g$ traverse leur bloc.
  \item \code{kill} : la réécriture propre (avec \code{own_value}, la valeur que
        \emph{l'appelant} attribue au site, voir \cref{sec:churn-point-fixe}) puis celles des
        \code{co} ; un slot écrit deux fois sur la chaîne prend le join et perd son
        \code{expr} (le bras relationnel n'a plus d'orthographe unique à comparer).
  \item \code{mine = let_bindings(site.cfg)} : les noms que le corps lie. Sous l'environnement
        de sortie du rendu, ils lisent $\Top$ (\code{shadowed}, \issue{162}) : le rendu peut
        lier le même nom --- un temporaire \code{__t0} surtout --- à autre chose, et une garde
        n'est jamais morte sur un nom qu'elle ne teste pas.
\end{itemize}
Le premier \code{dead_once_written} (\cref{chap:relations}, trois bras : valeur, relationnel,
membre) répond à la condition (a). S'il échoue, inutile d'aller plus loin.

\subsection{Seconde moitié : les revivers}

\rustsnippet[label={lst:churn-proof-2}]{src/engine/guards.rs}{198}{258}{\code{converges_under_all_writes}, seconde moitié}

\begin{itemize}
  \item \code{held} : les gardes du site qui \emph{tiennent} (\code{Invariance::holds}) ;
        ce sont les faits sous lesquels \emph{ce} site tire.
  \item \code{live(label)} : les revivers vivants du slot \code{label} --- ni le site, ni un
        \code{co}, du même slot qualifié, non convergents.
  \item Pour chaque slot $k$ de l'ensemble tueur : s'il est \code{foreign}, \code{false}.
        Pour chaque reviver \code{other} :
        \begin{enumerate}
          \item ses \code{facts} : ses gardes qui tiennent, et, si les deux corps diffèrent,
                dont \emph{aucun} nom n'est lié dans l'un ou l'autre corps
                (\code{names_unbound}) : un nom ne signifie la même chose dans deux corps
                que si tous deux le lisent dans la closure ;
          \item \code{contradicts(held, facts)} : un même conjoint (clé
                \code{call_free_key}) avec des polarités opposées --- l'autre site a tiré
                sous \code{!p}, \code{p} tient, donc le nôtre (sous \code{p}) ne tire jamais à
                côté de lui : ce reviver est \emph{sans objet}, on passe au suivant ;
          \item sinon on rétrécit l'environnement de sortie par les faits de \code{other},
                on réécrit : les slots de l'ensemble tueur \emph{autres} que $k$ gardent leur
                valeur tueuse \emph{s'ils n'ont aucun reviver vivant} (sinon ils lisent
                $\Top$, par absence de réécriture), et $k$ reçoit la valeur de \code{other}
                --- son \code{expr} n'étant lisible par le bras relationnel que si les corps
                sont les mêmes (\code{scope}) ;
          \item si les gardes ne sont pas mortes dans cet environnement, \code{false}.
        \end{enumerate}
\end{itemize}

\rustsnippet[label={lst:churn-contradicts}]{src/engine/guards.rs}{264}{273}{\code{contradicts}}

\begin{piege}
Le test de contradiction n'est sûr que parce qu'il porte sur des faits \emph{invariants}.
Deux gardes \code{s} et \code{!s} sur un état ne se contredisent pas d'un tour à l'autre :
c'est justement \code{s} qui bouge. \code{held} et \code{facts} sont filtrés par
\code{Invariance::holds} \emph{avant} \code{contradicts}, et un état n'y entre jamais.
\end{piege}

\subsection{Ce qui tient à travers la boucle : \code{Invariance}}

\rustsnippet[label={lst:churn-invariance}]{src/engine/guards.rs}{305}{315}{les faits d'invariance d'un composant}

\begin{definition}{Invariance à travers la boucle automatique}{churn-invariance}
La boucle re-rend sur les seules écritures d'état. \emph{Bouge} d'un tour à l'autre : un
état, ce qui en dérive (un memo, un callback, le résultat d'un hook), une allocation
fraîche, et tout nom qu'un corps du composant mute. \emph{Tient} : un littéral ; un nom lié
une fois (dans le corps ou dans le rendu) à quelque chose qui tient ; un nom que rien ne
lie --- un prop, un nom de module --- \emph{tant que les props tiennent}
(\code{props_hold}) ; un appel sur des entrées qui tiennent ; un résultat de hook dont le
résumé dit qu'il ne bouge que sur navigation (\code{SummaryValue::Held}), sauf si un corps
du programme navigue visiblement (\code{navigates}) ; un membre de littéral objet quand le
membre tient. Profondeur bornée à 8, au-delà : ne tient pas.
\end{definition}

La fonction \code{Invariance::holds} (\src{src/engine/guards.rs}{321}{390}) suit cette
définition cas par cas : \code{Lit} tient ; \code{Var} ne tient pas s'il est dans
\code{state_vals}, \code{memo_vals} ou \code{mutated}, sinon on suit sa liaison
(\code{binding}, corps puis rendu), et un nom que rien ne lie répond \code{props_hold} ;
\code{FieldAccess}/\code{IndexAccess} sont résolus à travers un littéral objet ou un
résultat de hook \og shaped\fg ; \code{BinOp}, \code{UnaryOp}, \code{Call} tiennent si
leurs entrées tiennent ; tout le reste (allocation, élément, autre hook) bouge.
\code{mutated} est l'union des racines mutées de tous les corps du composant
(\code{mutated_roots}) et des noms de module écrits par \emph{n'importe quel} composant
(\src{src/engine/churn.rs}{202}{224}). \code{navigates} est calculé une fois,
programme-entier, sur les rendus et tous les corps de hook \emph{sauf} les handlers
(\src{src/engine/churn.rs}{206}{216}) : une navigation depuis un handler exige un événement
utilisateur, elle n'est pas dans la boucle.

\begin{piege}
\og Un appel sur des entrées qui tiennent tient\fg{} (\srcline{src/engine/guards.rs}{294})
est une \emph{hypothèse}, pas une preuve : un callee opaque peut lire une horloge ou
naviguer. Elle est documentée dans \file{docs/limitations.md} et portée par \issue{161}.
C'est le seul point de la preuve où la soundness repose sur une hypothèse assumée plutôt
que sur une sur-approximation (\cref{sec:churn-limites}).
\end{piege}

\subsection{Retour sur les exemples}

\paragraph{\file{multi_writer.tsx} (\cref{ex:churn-multi-writer}).} Site
\code{setB({...})} de l'effet 2 : gardes $[(\code{b}, \code{false})]$, aucun \code{co}
(le corps n'a qu'une écriture), \code{kill} $= \{b \mapsto \text{réf. } \code{PerRender}\}$ :
morte, condition (a) satisfaite. \code{held} $= []$ (\code{b} est un état).
\code{live(b)} $= \{\code{setB(null)}\}$ de l'effet 4 : non gardé, \code{facts} $= []$ ;
\code{rewrites} $= \{b \mapsto \code{null}\}$ ; la garde \og \code{b} falsy\fg{} sur une
valeur \code{null} n'est pas $\Bot$ : \code{false}. L'arête $a \to b$ reste, le cycle est
rapporté.

\paragraph{\file{null_updater.tsx} (test \code{a_null_returning_updater_at_another_site_revives_the_guard}).}
Même schéma avec \code{if (s) setS(prev => null)} pour reviver : la sonde montre
\code{updater=Functional fresh=Not value.null=true}. Un updater fonctionnel est lu comme le
\emph{join de ses retours} (\cref{chap:relations}, \code{written::classify}), ici
\code{null} : la garde \code{!s} revit, l'arête \code{self_slot} reste, \Warning.

\paragraph{\file{settle_each_other.tsx} (\cref{ex:churn-settle}).} Site \code{setS} de
l'effet 2 : (a) morte sous $s \mapsto$ réf. \code{live(s)} $= \{\code{setS} \text{ de
l'effet 3}\}$ : sa garde \code{!s} ne tient pas (état), \code{facts} $= []$ ; réécriture
$s \mapsto$ réf. \code{PerRender} de l'autre site : \og \code{s} falsy\fg{} donne $\Bot$,
morte. \code{true}. Symétriquement pour l'effet 3. Les deux arêtes $b \to s$ sont tuées ;
seule $s \to b$ subsiste, sans cycle.

\subsection{Ce que la preuve garantit}

\begin{lemme}{Soundness du kill}{churn-kill-sound}
Supposons les hypothèses d'invariance de la \cref{def:churn-invariance} vraies du programme
concret, et supposons que tout site marqué \code{convergent} parmi les pairs tire un nombre
fini de fois dans la boucle automatique. Si \code{converges_under_all_writes} répond
\code{true} pour un site $p$, alors $p$ tire au plus une fois, puis un nombre fini de fois
supplémentaires au plus égal au nombre de tirs des sites convergents, dans toute boucle
automatique.
\end{lemme}
\begin{proof}[Esquisse]
Les gardes de $p$ sont conjonctives : il suffit qu'une soit fausse pour que $p$ ne tire pas.
Considérons un tour $t$ après le premier tir de $p$ et regardons, pour chaque slot $\ell$ de
l'ensemble tueur, la dernière écriture de $\ell$ avant $t$. Trois cas. (i) C'est $p$
lui-même ou un \code{co} : $\ell$ tient sa valeur tueuse, et les gardes sont mortes par (a).
(ii) C'est un reviver vivant $q$ : $q$ a tiré, donc ses faits invariants étaient vrais lors
de ce tir, et tiennent encore en $t$ ; si l'un contredit un fait de $p$, $p$ ne tire pas ;
sinon l'environnement de $t$ est inclus dans celui que la preuve a construit pour $q$
(faits de $q$, valeur de $q$ dans $\ell$, valeur tueuse ou $\Top$ ailleurs), où les gardes
sont mortes. (iii) C'est un site convergent : cela n'arrive qu'un nombre fini de fois, et
après chaque tel tir un tir de $p$ restaure la valeur tueuse. Un site étranger n'écrit
jamais un slot tueur (\code{foreign} exclu). Donc $p$ ne tire, hors les cas (iii), plus
jamais.
\end{proof}

Le lemme est conditionnel, et c'est voulu : la seconde hypothèse est \emph{discharged} par
le plus petit point fixe de la section suivante, la première reste une hypothèse du
projet.

% ---------------------------------------------------------------------------
\section{Revivers et plus petit point fixe des sites convergents}
\label{sec:churn-point-fixe}

\subsection{Un reviver qui tire une fois par requête}

La preuve de la section précédente traite un reviver \emph{vivant} comme éternel : s'il peut
ranimer la garde une fois, il pourra la ranimer toujours. C'est faux pour un reviver qui,
lui aussi, converge. Le motif vient du dépôt \code{twenty} (\issue{160}).

\begin{exemple}{Un drapeau de requête et sa continuation}{churn-reviver}
\begin{lstlisting}[language=TypeScript,caption={\file{/tmp/ch15/reviver.tsx}},label={lst:churn-reviver}]
import { useState, useEffect } from 'react';
declare function refetch(): Promise<void>;
export function Reviver({ version }: { version: { id: string } }) {
  const [req, setReq] = useState(false);
  const [seeded, setSeeded] = useState<string | undefined>(undefined);
  useEffect(() => {
    if (!req) return;
    setReq(false);
    void refetch().then(() => { setSeeded(undefined); });
  }, [req]);
  useEffect(() => {
    if (seeded === version.id) return;
    setSeeded(version.id);
  }, [version, seeded]);
  return <div>{seeded}</div>;
}
\end{lstlisting}
\begin{sortie}
  Reviver  (4 hooks)  /tmp/ch15/reviver.tsx  ✓
    verified  infinite-loop  no effect diverges into an infinite render loop
\end{sortie}
Test : \code{a_reviver_that_resets_its_own_flag_}\allowbreak\code{fires_once_per_request}.
\end{exemple}

Le second effet s'éteint \emph{relationnellement} : la garde compare \code{seeded} à
\code{version.id}, l'écriture pose \code{version.id} dans \code{seeded}, les deux côtés sont
égaux au rendu suivant. Mais la continuation \code{.then} du premier effet remet
\code{undefined} dans \code{seeded} : un reviver. Il n'est pourtant planifié que si
\code{req} est vrai, et la \emph{même passe} remet \code{req} à \code{false} : il tire une
fois par requête externe, la paire converge en deux tours. Pour le voir, il faut trois
choses.

\begin{sortie}
  slot=0 setter=setReq    line=Some(8)  region=Effect(2) phase=Effect   block=Some(3) guard_block=Some(3) fresh=Not
  slot=1 setter=setSeeded line=Some(9)  region=Effect(2) phase=Deferred block=None    guard_block=Some(3) fresh=Not
  slot=1 setter=setSeeded line=Some(13) region=Effect(3) phase=Effect   block=Some(3) guard_block=Some(3) fresh=Maybe value.ref=Unknown
-- effect_triggers
  hook=2 dep=0 slot=(Reviver, 0) exact=true
  hook=3 dep=1 slot=(Reviver, 1) exact=true
-- registrations
  effect=2 display=.then registrar=then firing=Once timing=Deferred block_id=Some(3) line=Some(9) pairing=Unpaired
=== churn graph
\end{sortie}

\begin{enumerate}
  \item \textbf{Une écriture différée a des gardes.} La ligne \code{setSeeded(undefined)}
        est \code{Deferred} (\code{block=None}) mais porte \code{guard_block=Some(3)} : le bloc
        de l'instruction \code{refetch().then(...)} qui l'a planifiée. La continuation ne
        tourne que si la passe qui l'a planifiée a pris toutes les branches au-dessus ;
        ses gardes sont donc celles du bloc 3 : \code{guard_chain(3)} $= [3, 2, 0]$,
        \code{site_guards} $= [(\code{req}, \code{true})]$ (la \cref{fig:churn-kill-set}).
  \item \textbf{L'ensemble tueur contient \code{setReq(false)}.} Cette écriture est
        synchrone, locale, dans le bloc 3 qui est sur la chaîne : c'est un \code{co}. Sous
        $\code{req} \mapsto \code{false}$, la garde $(\code{req}, \code{true})$ est $\Bot$ :
        le reviver \emph{converge}.
  \item \textbf{Un reviver convergent n'est pas vivant.} Une fois \code{setSeeded(undefined)}
        prouvé, il sort de \code{live(seeded)} pour le site \code{setSeeded(version.id)}, qui
        converge à son tour.
\end{enumerate}

\begin{figure}[htbp]\centering
\begin{tikzpicture}[node distance=7mm and 16mm]
  \node[cfgbloc,fill=ra-bg] (b0) {bloc 0\\Branch \{ cond: !req \}};
  \node[cfgbloc,below left=of b0] (b1) {bloc 1\\Return};
  \node[cfgbloc,fill=ra-bg,below right=of b0] (b2) {bloc 2\\Jump(3)};
  \node[cfgbloc,fill=ra-bg,draw=ra-blue,very thick,below=of b2] (b3) {bloc 3\\setReq(false)\\refetch().then(() => \{ setSeeded(undefined) \})};
  \draw[fleche] (b0) -- node[etiquette,left] {IfTrue} (b1);
  \draw[fleche] (b0) -- node[etiquette,right] {IfFalse} (b2);
  \draw[fleche] (b2) -- (b3);
  \coordinate (mid) at ($(b1.south)!0.5!(b3.south)$);
  \coordinate (anch) at (mid |- b3.south);
  \node[etiquette,text width=112mm,align=left,below=5mm of anch,anchor=north] {%
    \code{guard_chain(cfg, 3) = [3, 2, 0]} (blocs grisés) ;\\
    \code{site_guards(cfg, 3) = [(req, true)]} : l'arête $0 \to 2$ est \code{IfFalse} sur \code{!req}, dépliée en \code{req} vrai ;\\
    ensemble tueur du site \code{setSeeded(undefined)} : $\{\code{seeded} \mapsto \code{undefined},\ \code{req} \mapsto \code{false}\}$ ---
    \code{setReq(false)} est un \code{co} (même corps, slot local, \code{block = Some(3)} sur la chaîne).};
\end{tikzpicture}
\caption{Le corps du premier effet de l'\cref{ex:churn-reviver} après lowering (blocs de la sonde IR), la chaîne de gardes du bloc 3 et l'ensemble tueur de l'écriture différée. Le site différé n'a pas de \code{block} mais hérite du \code{guard_block} de l'instruction qui le planifie.}
\label{fig:churn-kill-set}
\end{figure}

\subsection{Le point fixe des sites convergents}

\rustsnippet[label={lst:churn-fixpoint}]{src/engine/churn.rs}{450}{485}{le plus petit point fixe des sites convergents}

Le vecteur \code{convergent} part de \og aucun\fg. À chaque tour, pour chaque composant, on
projette les sites en \code{peers} \emph{avec le drapeau courant}, puis on tente la preuve sur
chaque site local non encore prouvé, avec sa \emph{valeur entière} (\code{peers[k].value}).
On s'arrête quand un tour ne marque plus rien. Les lignes étrangères
(\code{owner.is_some()}) ne sont jamais tentées.

\begin{theoreme}{Le plus petit point fixe des sites convergents}{churn-lfp}
Soit $S$ l'ensemble des sites d'un programme et, pour $X \subseteq S$, $F(X)$ l'ensemble des
sites $p$ tels que \code{converges_under_all_writes} répond \code{true} quand les pairs
marqués convergents sont exactement ceux de $X$. Alors $F$ est monotone
($X \subseteq Y \Rightarrow F(X) \subseteq F(Y)$), l'itération $X_0 = \emptyset$,
$X_{n+1} = X_n \cup F(X_n)$ stationne en au plus $|S| + 1$ tours sur
$\lfp(X \mapsto X \cup F(X))$, et le résultat ne dépend pas de l'ordre de visite des
composants ni des sites.
\end{theoreme}
\begin{proof}
Marquer un site convergent ne fait que \emph{retirer} des éléments de \code{live(label)}
pour les autres : moins de revivers à réfuter, chaque condition de la preuve est
monotone décroissante en \code{live}, donc $F$ est monotone. La suite $X_n$ est croissante
dans un ensemble fini, d'où la terminaison ; un tour utile marque au moins un site, d'où la
borne. L'indépendance de l'ordre est celle de tout plus petit point fixe d'un opérateur
monotone sur un treillis fini (\cref{def:point-fixe}) ; le code le note à
\src{src/engine/churn.rs}{446}{449}.
\end{proof}

\begin{remarque}
L'itération est \emph{à la Jacobi} : \code{peers} est construit une fois au début de chaque
tour à partir de \code{convergent}, avant la boucle interne. Un site prouvé au tour $n$
n'est vu convergent par ses pairs qu'au tour $n+1$. Le résultat n'en dépend pas (théorème),
mais un déroulé à la main doit compter les tours ainsi.
\end{remarque}

\begin{table}[htbp]\centering\footnotesize
\begin{tabularx}{\linewidth}{>{\raggedright\arraybackslash}p{36mm}>{\raggedright\arraybackslash}p{30mm}>{\raggedright\arraybackslash}X>{\raggedright\arraybackslash}X}
\toprule
Site & Gardes & Tour 1 & Tour 2 \\
\midrule
\code{setReq(false)} (effet 2, bloc 3) & $(\code{req}, \top)$ & \checkmark{} (a) : $\code{req} \mapsto \code{false}$ tue ; aucun reviver de \code{req} & --- \\
\addlinespace
\code{setSeeded(undefined)} (effet 2, \code{.then}) & $(\code{req}, \top)$ & \checkmark{} (a) : \code{co} $= \{\code{setReq}\}$ tue ; (b) : le reviver \code{setSeeded(version.id)} réécrit \code{seeded}, \code{req} garde \code{false} & --- \\
\addlinespace
\code{setSeeded(version.id)} (effet 3, bloc 3) & $(\code{seeded} \equiv$ \newline $\code{version.id}, \bot)$ & $\times$ (a) tient ; (b) : le reviver \code{setSeeded(undefined)} est encore vivant dans le snapshot & \checkmark{} \code{live(seeded)} $= \emptyset$ \\
\midrule
\code{changed} & & \code{true} & \code{true} ; tour 3 : \code{false}, sortie \\
\bottomrule
\end{tabularx}
\caption{Les tours du point fixe des sites sur l'\cref{ex:churn-reviver}. $(c, \top)$ note une garde $c$ prise vraie, $(c, \bot)$ prise fausse ; (a) et (b) renvoient aux deux conditions de la \cref{def:churn-kill-set}. Le déroulé a été observé en instrumentant une copie de \code{ChurnGraph::build} (une trace par site et par tour) : les deux premiers sites sont marqués au tour 1, le troisième au tour 2, et le tour 3 ne change rien. Sur \file{reviver_unreset.tsx}, le même instrument montre un seul tour, sans aucun site marqué.}
\label{tab:churn-iterations}
\end{table}

Détaillons la case $\times$ du tour 1, la seule subtile. Le site \code{setSeeded(version.id)}
satisfait (a) par le bras relationnel : la garde est
\code{BinOp { Eq, seeded, version.id }} prise fausse, l'écriture pose \emph{verbatim}
\code{version.id} dans \code{seeded}, ni l'un ni l'autre n'est une orthographe fraîche
(\code{version.id} est un membre de prop) : après l'écriture l'égalité est vraie, la
branche fausse est morte. Pour (b), \code{live(seeded)} contient \code{setSeeded(undefined)}
tant qu'il n'est pas marqué. Ses faits : $(\code{req}, \top)$, mais \code{req} est un état,
il ne tient pas, \code{facts} $= []$. On réécrit $\code{seeded} \mapsto \code{undefined}$
avec \code{scope = None} (autre corps) : le bras relationnel est inapplicable, le bras
membre aussi (le chemin est le slot nu), et le bras valeur rétrécit
\code{seeded === version.id} --- dont le côté droit n'est pas un littéral : l'environnement
est inchangé, rien n'est $\Bot$. Échec. Au tour 2, le snapshot marque le reviver
convergent : \code{live(seeded)} est vide, la boucle sur les revivers ne tourne pas,
\code{true}.

\subsection{Pourquoi pas le plus grand point fixe}

On pourrait partir de \og tous convergents\fg{} et retirer les sites que la preuve réfute.
Ce plus grand point fixe est \emph{insound}, et le contre-exemple est court.

\begin{exemple}{Deux écritures qui se raniment l'une l'autre}{churn-mutual}
\begin{lstlisting}[language=TypeScript,caption={\file{/tmp/ch15/null_updater.tsx}},label={lst:churn-null-updater}]
import { useState, useEffect } from 'react';
export function NullUpdater() {
  const [s, setS] = useState(null);
  useEffect(() => { if (!s) setS({ fresh: true }); }, [s]);
  useEffect(() => { if (s) setS(prev => null); }, [s]);
  return <div>{s ? 'y' : 'n'}</div>;
}
\end{lstlisting}
\begin{sortie}
  NullUpdater  (3 hooks)  null_updater.tsx
    warn   infinite-loop  [hook:0]  (line 4:2)  this effect may store a fresh reference into state `s` which its deps react to: possible infinite render loop
       → a fresh value is written to state `s` here [hook:1] (line 4:28)
\end{sortie}
\end{exemple}

Appelons $A$ le site \code{setS({...})} (garde \code{!s}) et $B$ le site
\code{setS(prev => null)} (garde \code{s}). Chacun meurt sous sa propre écriture : $A$ pose
une référence truthy, $B$ pose \code{null}. Chacun est ranimé par l'autre : $B$ remet
\code{s} falsy, $A$ le remet truthy. \emph{Si l'on suppose $B$ convergent}, $A$ n'a plus de
reviver vivant et est prouvé ; \emph{si l'on suppose $A$ convergent}, $B$ est prouvé. Le
plus grand point fixe contient donc $\{A, B\}$ : les deux arêtes seraient tuées et le
programme --- qui boucle \code{{...}} $\to$ \code{null} $\to$ \code{{...}} $\to$ … dès le
montage --- serait silencieux. Le plus petit point fixe ne marque ni l'un ni l'autre : au
tour 1, chacun voit l'autre vivant et échoue ; rien ne change ; sortie. L'arête
\code{self_slot} $s \to s$ reste et le bras self-churn émet le \Warning{} observé.

\begin{theoreme}{Soundness du plus petit point fixe}{churn-lfp-sound}
Sous les hypothèses d'invariance, tout site de $\lfp$ tire un nombre fini de fois dans
toute boucle automatique. Le plus grand point fixe de $F$ n'a pas cette propriété.
\end{theoreme}
\begin{proof}
Par récurrence sur le tour $n$ où un site $p$ entre dans $X_n$. Les pairs que la preuve
de $p$ tient pour convergents sont dans $X_{n-1}$ et, par hypothèse de récurrence, tirent
un nombre fini de fois ; le \cref{lem:churn-kill-sound} conclut pour $p$. Pour la
seconde partie, l'\cref{ex:churn-mutual} exhibe $\{A, B\} \subseteq \mathrm{gfp}$ avec une
boucle infinie.
\end{proof}

\begin{decision}[ADR-42 \S6, amendement \issue{160}]
Le point fixe \emph{minimal} est un choix argumenté en une phrase dans le code même :
\og a greatest fixpoint would read two sites that revive each other as convergent, which
is exactly a loop\fg{} (\src{src/engine/guards.rs}{41}{42}). L'alternative --- une colonne
\code{settles} sur la ligne, calculée une fois --- avait déjà été refusée à la tranche 4 :
deux appelants demandent deux valeurs (voir ci-dessous), et désormais la réponse dépend
aussi de l'état des autres sites, ce qu'une colonne ne peut porter. Mesure du commit
\snapshotCommit{} (\file{docs/precision-log.md}) : 1\,499 $\to$ 1\,498 sur le corpus, trois
lignes retirées et deux ajoutées, dont deux sont une même ligne reformulée.
\end{decision}

Le programme frère \file{reviver_unreset.tsx} (test
\code{a_reviver_that_keeps_its_flag_still_revives}), où \code{setReq(false)} est retiré et
\code{req} vaut \code{true} pour toujours, donne un \Warning{} sur \code{seeded} : le
reviver différé n'a plus de \code{co} qui tue sa garde $(\code{req}, \top)$, il reste vivant,
et \code{setSeeded(version.id)} n'est jamais prouvé. La preuve ne lit pas une garde comme
mourant sous une écriture qui n'a pas lieu.

\subsection{Le kill des arêtes : la partie référence}

Le point fixe des sites fait, l'étape des arêtes réutilise la même preuve, mais avec une
autre valeur propre.

\rustsnippet[label={lst:churn-killed}]{src/engine/churn.rs}{502}{529}{le kill d'une arête}

\rustsnippet[label={lst:churn-reference-part}]{src/engine/written.rs}{202}{209}{\code{reference_part}}

Une arête affirme un churn de \emph{référence} : ce qu'elle transporte d'un rendu à l'autre
est la partie référence de la valeur écrite. Le site est donc lu, pour sa propre écriture,
comme \code{reference_part(w.written.value)} --- les parties primitives, qui ne peuvent
pas échouer \code{Object.is} fraîchement, sont abandonnées, ce qui évite de perdre la
preuve à un $\Top$ résiduel ; une partie référence $\Bot$ donne $\Bot$ (aucune référence
ne peut être stockée, le churn prétendu est vide). Les \emph{autres} sites, eux, raniment avec
tout ce qu'ils stockent, \code{null} compris : dans la preuve, \code{other.value} est la
valeur entière. Le point fixe des sites, à l'inverse, lit la valeur entière du site propre
(\code{peers[k].value}) : il répond à la question \og ce site tire-t-il une fois ?\fg{} et
non \og cette arête est-elle un churn ?\fg. Deux appelants, deux valeurs : c'est
l'argument d'\adr{42} \S6 (tranche 4) pour une fonction plutôt qu'une colonne.

\begin{piege}
Le \code{killed} de l'\cref{lst:churn-killed} est une disjonction :
\code{convergent[idxs[k]] || converges_under_all_writes(..., &own_value, ...)}. Un site que le
point fixe n'a pas marqué (parce que sa valeur \emph{entière}, \code{null} compris, ne tue
pas ses gardes) peut encore voir son \emph{arête} tuée, parce que la partie référence de sa
valeur, elle, les tue. Le premier disjoint, lui, est un raccourci : un site déjà prouvé
convergent n'a pas besoin d'une seconde preuve, puisque ne pas tirer implique ne pas
churner. Ne pas confondre les deux questions en lisant une trace : \og site non
convergent\fg{} n'implique pas \og arête présente\fg.
\end{piege}

% ---------------------------------------------------------------------------
\section[Plusieurs composants]{Plusieurs composants : lignes étrangères, \code{props_hold}, plafond Warning}
\label{sec:churn-cross}

\subsection{Un setter passé en prop}

\begin{exemple}{Un enfant qui réécrit l'état de son parent}{churn-cross}
\begin{lstlisting}[language=TypeScript,caption={\file{/tmp/ch15/cross.tsx}},label={lst:churn-cross}]
import { useState, useEffect } from 'react';
export function Parent() {
  const [data, setData] = useState({ n: 0 });
  return <Child value={data} onUpdate={setData} />;
}
function Child({ value, onUpdate }) {
  useEffect(() => { onUpdate({ n: value.n, seen: true }); }, [value]);
  return <div/>;
}
\end{lstlisting}
\begin{sortie}
  Child  (1 hooks)  /tmp/ch15/cross.tsx
    warn   cross-component-infinite-loop  [hook:0]  (line 7:2)  this effect calls `onUpdate`, a state setter of parent `Parent` (its deps do not provably gate it, so the effect can re-run every render). Parent re-renders → child re-renders → effect fires again: infinite loop
  Parent  (1 hooks)  /tmp/ch15/cross.tsx  ✓
\end{sortie}
\begin{sortie}
=== component Child (ComponentId(0))
  slot=0 owner=Some("Parent") setter=onUpdate line=Some(7) region=Effect(0) phase=Effect block=Some(0) fresh=Fresh value.ref=PerRender
-- effect_triggers
  hook=0 dep=0 slot=(Parent, 0) exact=false
=== churn graph
  edge#0 (Parent, 0) -> (Parent, 0) strength=May carrier=Child effect=0 line=Some(7) no_deps=false self_slot=false
  cycle edges=[0] all_must=false cross_component=true
\end{sortie}
\end{exemple}

Dans l'environnement de \code{Child}, \code{onUpdate} vaut
\code{ComponentSetter(Parent, 0)} (passe descendante, \cref{chap:inter-composants}) : la
ligne d'écriture est \emph{étrangère}, \code{owner=Some("Parent")}, et \code{node_of} la
place sur le nœud \code{(Parent, 0)}. Le dep \code{value} est un prop dont la valeur est
\code{Versioned({(Parent, 0)})} : le déclencheur est \code{exact=false}. Trois
conséquences s'enchaînent :
\begin{enumerate}
  \item l'arête est \code{May} (\code{versioned}, \srcline{src/engine/churn.rs}{592}) ;
  \item elle n'est pas \code{self_slot} : la condition \code{x == node && x.0 == f.comp}
        (\srcline{src/engine/churn.rs}{589}) échoue car le slot appartient à \code{Parent}
        et l'effet à \code{Child} ; elle entre donc dans la recherche de cycles, où elle
        forme un cycle de longueur 1 ;
  \item \code{make_cycle} compte deux composants (propriétaire du nœud et porteur) :
        \code{cross_component = true}.
\end{enumerate}

\begin{decision}[ADR-18, plafond Warning]
Un cycle qui traverse deux composants est plafonné à \Warning{} et nommé
\regle{cross-component-infinite-loop}. La raison est structurelle : le must-rerun
à travers un prop est \emph{indémontrable} au commit \snapshotCommit, parce qu'un dep
qui est un prop est \code{Versioned}, jamais le slot exact --- le composant enfant ne voit
pas \emph{quel} champ du parent a bougé. \adr{18} note la voie d'amélioration (une
provenance de prop entier au site d'appel JSX) sans la prendre.
\end{decision}

Le diagnostic observé ci-dessus n'est pourtant pas celui du bras graphe : c'est le
\emph{bras cross-component historique} de \regle{infinite-loop} qui parle (\og this effect
calls \code{`onUpdate`}, a state setter of parent\fg). L'effet étant déjà dans
\code{reported_effects}, le bras graphe le saute
(\src{src/rules/impls/infinite_loop.rs}{277}{279}) : un rapport par effet et par classe de
panne.

\subsection{Quand les props bougent : \code{props_hold}}

La preuve de convergence lit des faits sur les props (\og \code{p} tient\fg) : un prop n'est
pas un état du composant. Mais un prop \emph{est} un état du parent, et si la boucle passe
par le parent, il bouge.

\begin{exemple}{Un cycle qui entre par les props}{churn-props-move}
\begin{lstlisting}[language=TypeScript,caption={\file{/tmp/ch15/props_move.tsx}},label={lst:churn-props-move}]
import { useState, useEffect } from 'react';
function Child({ p, onChange }: { p: any; onChange: (v: any) => void }) {
  const [x, setX] = useState(null);
  const [s, setS] = useState(null);
  useEffect(() => { setX({ p }); }, [p]);
  useEffect(() => { if (p && !s) setS({ v: 1 }); }, [x]);
  useEffect(() => { if (!p) setS(null); }, [x]);
  useEffect(() => { onChange(s ? null : { k: 1 }); }, [s]);
  return null;
}
export function Parent() {
  const [p, setP] = useState<any>({ k: 1 });
  return <Child p={p} onChange={setP} />;
}
\end{lstlisting}
\begin{sortie}
  Child  (6 hooks)  props_move.tsx
    warn   cross-component-infinite-loop  [hook:2]  (line 5:2)  these effects may form a state-update cycle (`x` → `s` → `p` of `Parent` → `x`) where each step may store a fresh reference that re-runs the next effect: possible infinite render loop
       → a fresh value is written to state `x` here [hook:2] (line 5:20)
       → cycle continues: this effect freshly stores state `s` [hook:3] (line 6:33)
       → cycle continues: this effect freshly stores state `p` of `Parent` [hook:5] (line 8:20)
    warn   cross-component-infinite-loop  [hook:3]  (line 6:2)  these effects may form a state-update cycle (`x` → `s` → `p` of `Parent` → `x`) …
    warn   cross-component-infinite-loop  [hook:5]  (line 8:2)  this effect calls `onChange`, a state setter of parent `Parent` (its deps do not provably gate it, so the effect can re-run every render). …
\end{sortie}
Test : \code{props_move_on_a_cycle_}\allowbreak\code{that_enters_through_another_edge}.
\end{exemple}

Le cycle est $(\code{Parent}, p) \to x \to s \to (\code{Parent}, p)$ : il \emph{entre} dans
\code{Child} par le prop \code{p} (arête versionnée vers \code{x}) et en \emph{sort} par le
setter-prop \code{onChange}. L'arête $x \to s$ est locale, et les deux sites de \code{s}
tirent sous \code{p} et \code{!p} : si ces faits \emph{tenaient}, ils se contrediraient et
la garde \code{!s} mourrait --- l'arête serait tuée et le cycle perdu. Or \code{p} bouge le
long du cycle. La \cref{fig:churn-props-move} dessine le graphe.

\begin{figure}[htbp]\centering
\begin{tikzpicture}[node distance=12mm and 20mm]
  \node[noeud] (x) at (0,0) {$x$};
  \node[noeud,right=of x] (s) {$s$};
  \node[noeud] (p) at (2.2,-2.2) {$p$};
  \begin{scope}[on background layer]
    \node[draw=ra-gray,dashed,rounded corners=4pt,fit={(x) (s)},inner sep=6mm,label={[etiquette]above left:\code{Child}}] {};
    \node[draw=ra-gray,dashed,rounded corners=4pt,fit={(p)},inner sep=5mm,label={[etiquette]below:\code{Parent}}] {};
  \end{scope}
  \draw[flechemay] (x) -- node[etiquette,above] {\code{e3}, sous \code{p && !s}} (s);
  \draw[flechemay] (s) -- node[etiquette,right,xshift=1mm] {\code{e5}, via \code{onChange}} (p);
  \draw[flechemay] (p) -- node[etiquette,left,xshift=-1mm] {\code{e2}, dep \code{p} versionné} (x);
\end{tikzpicture}
\caption{Le graphe de \file{props_move.tsx} : un cycle de trois arêtes \code{May} dont un nœud appartient à \code{Parent} (\code{cross_component}). \code{Child} porte les trois arêtes ; comme un de ses déclencheurs nomme un slot du parent, \code{props_hold} est faux et les faits sur \code{p} ne sont pas lus.}
\label{fig:churn-props-move}
\end{figure}

La solution est un fait par composant, calculé avant toute preuve
(\srcline{src/engine/churn.rs}{232}) :
\code{props_hold = effect_triggers.iter().all(|t| t.slot.0 == comp)}. Aucun effet du
composant ne réagit au slot d'un autre composant $\Rightarrow$ aucune boucle n'y entre par
ses props $\Rightarrow$ ses props tiennent le long de toute boucle qui passe par ses arêtes.
Un seul dep versionné par un slot du parent, et les props peuvent bouger sur \emph{tout}
cycle qui traverse le composant, quelle que soit l'arête empruntée : les preuves ne lisent
alors plus aucun fait sur eux (\code{Invariance::holds} rend \code{props_hold} pour un
nom que rien ne lie, \srcline{src/engine/guards.rs}{353}). Ici \code{Child} a un
déclencheur sur \code{(Parent, 0)}, donc \code{props_hold = false}, les faits \code{p} et
\code{!p} sont écartés, aucune contradiction, l'arête $x \to s$ survit, le cycle est
rapporté.

\begin{remarque}
Le programme reçoit \emph{trois} \Warning{} pour une seule boucle : deux du bras graphe
(effets 2 et 3, porteurs d'arêtes) et un du bras cross-component historique (effet 5, qui
appelle le setter du parent). La déduplication de \srcline{src/rules/impls/infinite_loop.rs}{277}
ne joue que par effet : l'effet 5 est sauté par le bras graphe parce que déjà rapporté, les
effets 2 et 3 ne le sont pas. C'est un choix de \adr{18} (\og un rapport par effet
porteur\fg), pas un défaut.
\end{remarque}

\subsection{Ce qui ne remonte pas}

Le graphe ne voit une écriture étrangère que si le parent a été analysé \emph{descendant}
et que \code{ComponentSetter} a coulé jusqu'à l'enfant. Un parent analysé seulement en intra
laisse l'enfant sans ligne étrangère, donc sans arête et sans
\regle{cross-component-infinite-loop} (\issue{20}, \code{precision-fn}). Les Consequences
d'\adr{42} l'enregistrent, en reprenant \adr{29}, comme affaiblissement connu : \og rows
exist only for cycles the graph sees; prop-mediated edges need the parent slot flowed
top-down (\#20)\fg.

% ---------------------------------------------------------------------------
\section{De l'arête au diagnostic : la frontière de certification}
\label{sec:churn-certification}

Le graphe dit \code{Must} ou \code{May} ; il ne \emph{frappe} rien. La sévérité est
décidée côté règles, et jamais en faisant confiance à un booléen fourni par le graphe :
elle est \emph{re-dérivée} au point de connaissance (\cref{chap:api-regles},
\cref{def:certified}).

\subsection{Le bras graphe et \code{must_effect_cycle}}

\rustsnippet[label={lst:churn-arm3}]{src/rules/impls/infinite_loop.rs}{281}{292}{le bras multi-effets choisit la règle et demande la preuve}

\rustsnippet[label={lst:churn-must-effect-cycle}]{src/rules/api/query.rs}{1002}{1031}{\code{must_effect_cycle} re-dérive les deux faits}

\code{must_effect_cycle} recalcule depuis les arêtes brutes ce que \code{ChurnCycle} porte
déjà : toutes les arêtes \code{Must} \emph{et} un seul composant impliqué. Si les deux
tiennent, il \emph{frappe} (\code{Certified::mint}) une preuve \code{EffectCycleProof} dont
la provenance est la première arête du cycle --- l'écriture qui relance l'effet porteur.
Sinon \code{MustResult::None}, et le diagnostic est un \Warning. La règle choisit le nom
selon \code{cycle.cross_component} et la sévérité selon la \emph{preuve} : un cycle
cross-component tout-Must est possible (deux arêtes \code{Must} locales et une auto-arête
étrangère, par exemple) et reste un \Warning, et sa formulation suit la preuve (\og may
form\fg), pas \code{all_must} (\src{src/rules/impls/infinite_loop.rs}{294}{296}).

\begin{table}[htbp]\centering\small
\begin{tabularx}{\linewidth}{cccX}
\toprule
$|$cycle$|$ & \code{no_deps} & preuve & message (début) \\
\midrule
1 & oui & oui & \og this effect has no dependency array and stores a fresh reference … infinite render loop\fg \\
1 & oui & non & \og … and may store a fresh reference … possible infinite render loop\fg \\
1 & non & oui & \og this effect stores a fresh reference into state … which its own deps react to\fg \\
1 & non & non & \og this effect may store a fresh reference … which its own deps react to\fg \\
$\geq 2$ & --- & oui & \og these effects form a state-update cycle (\emph{chemin}) …\fg \\
$\geq 2$ & --- & non & \og these effects may form a state-update cycle (\emph{chemin}) …\fg \\
\bottomrule
\end{tabularx}
\caption{Les six formulations du bras graphe (\src{src/rules/impls/infinite_loop.rs}{298}{334}). La colonne \og preuve\fg{} est \code{cycle_proof.is_some()}, jamais \code{all_must}.}
\label{tab:churn-messages}
\end{table}

Le témoin (\cref{def:witness}) est une étape \code{Write} (valeur \code{Fresh}) au
\code{write_span} de l'arête portée, puis une étape \code{CycleEdge} par autre arête du cycle
portée par le même composant : c'est la liste \og cycle continues\fg{} des sorties de ce
chapitre. Le chemin \code{`a`}~→ \code{`b`}~→ \code{`a`} est rendu par \code{cycle_path}
(\src{src/rules/helpers/cycles.rs}{54}{70}), qui nomme un slot étranger
\code{`p` of `Parent`} (\code{node_display}) ; ces deux fonctions vivent côté règles
\og parce qu'un nom d'affichage est un rendu, pas un fait\fg{} (doc de module,
\src{src/rules/helpers/cycles.rs}{1}{8}).

\subsection{Le bras self-churn re-dérive must-reach}

\rustsnippet[label={lst:churn-arm4}]{src/rules/impls/infinite_loop.rs}{441}{460}{le bras self-churn ne fait pas confiance à \code{e.strength}}

Le bras self-churn (\cref{chap:infinite-loop}) lit sa partition --- les arêtes du composant,
de l'effet, non \code{no_deps}, de slot local vers slot local
(\src{src/rules/impls/infinite_loop.rs}{427}{433}) --- et, pour une arête \code{self_slot},
\emph{recalcule} l'ensemble des blocs d'écritures fraîches depuis \relation{slot_writers} et
demande \code{must_on_all_paths}. \code{e.strength == Must} est une condition nécessaire
(elle évite un calcul inutile), jamais suffisante : la preuve est frappée par la
must-primitive. Les arêtes non \code{self_slot} de la partition alimentent l'\Info{}
\og freshly recreates state … outside its deps\fg, supprimée pour les écritures
\code{covered} par un cycle rapporté (\cref{ex:churn-dag}).

\subsection{L'ancre \code{churn_cycles} des règles déclaratives}

\rustsnippet[label={lst:churn-cyclerow}]{src/rules/helpers/cycles.rs}{72}{93}{une ligne \code{churn_cycles}}

\begin{decision}[ADR-29]
Les règles déclaratives (\cref{chap:regles-declaratives}) sont \emph{mono-ancre} : une règle
parle d'un composant. Le graphe étant programme-entier, \adr{29} le \emph{projette} : une
ligne par arête de cycle \emph{portée par un effet de ce composant}
(\code{collect_cycle_rows}, \src{src/rules/helpers/cycles.rs}{101}{135}), avec le chemin
rendu, deux booléens exacts (\code{cross_component}, \code{all_must} --- le may vit dans le
graphe, pas dans la ligne) et le span de l'écriture comme identité (\adr{24}). Une arête
sans span ne donne \emph{aucune} ligne : un canal de findings manqués enregistré, jamais
un finding pointant nulle part. Refusés : un statut \og couvert nativement, ne compte
pas\fg{} ; fusionner les deux bras (\adr{20} item 2). L'ancre est plafonnée à \Warning{}
structurellement : seule la règle native atteint \Error, via \code{must_effect_cycle}.
\end{decision}

\subsection{Construit une fois}

\code{ChurnGraph::build} lit tous les composants ; il est appelé une fois par programme,
paresseusement, par l'accesseur \code{ProgramRelations::churn} : un \code{OnceLock} rempli
à la première demande par \code{get_or_init}
(\src{src/engine/program_relations.rs}{39}{40}, \cref{chap:relations}). Il est ensuite
partagé par toutes les passes de règles. Un test le garantit par le compteur \code{BUILDS} :
\code{churn_graph_is_built_once_per_program}
(\src{src/rules/impls/infinite_loop.rs}{631}{684}). C'est la correction de \issue{86}, où le
graphe était reconstruit par composant et la phase des règles devenait quadratique en nombre
de composants.

% ---------------------------------------------------------------------------
\section{Limites connues au commit \snapshotCommit}
\label{sec:churn-limites}

\subsection{Faux négatifs connus}

\begin{description}
  \item[\issue{157}, valeur dérivée du slot écrit.] \code{const d = useMemo(() => ({...s}), [s]);
        useEffect(() => setS(d), [d])} : la valeur écrite est \code{Versioned({s})}, que
        \code{written::classify} lit \code{Not} frais (lecture d'\adr{17}) ; aucune arête,
        boucle silencieuse. Le correctif naïf (\code{Versioned} $\Rightarrow$ \code{Maybe})
        transforme le motif \og mirror state\fg{} --- recopier un prop dans un état --- en
        \Warning{} massif ; l'issue reste ouverte, étiquetée \code{soundness-bug}.
  \item[\issue{161}, l'hypothèse d'invariance.] Un appel sur des entrées qui tiennent est
        supposé rendre la même valeur ; une navigation cachée dans un callee opaque ou dans
        un objet routeur passé en prop n'est pas vue (\file{docs/limitations.md}).
  \item[\issue{20}, parent analysé en intra.] Voir \cref{sec:churn-cross}.
\end{description}

\subsection{Faux positifs assumés (Warning au plus)}

Un slot convergent que \emph{un autre composant} écrit aussi (résidu de \issue{39} :
\code{foreign} n'est jamais tué) ; une garde sur le résultat d'un hook non inliné ou d'un
store externe, lu comme bougeant (\issue{161}) ; un setter-prop typé primitif dont
l'argument $\Top$ est lu comme référence possible (\issue{159}) ; une garde disjonctive ou
arithmétique sur la valeur comparée (\issue{91}) ; un enfant rendu depuis un \code{.map} ou
atteint via un composant intermédiaire, lu comme remontant (\issue{162}, côté fail-closed).

\subsection{Granularité et forme du résultat}

Le graphe est granulaire \emph{au slot} : savoir si l'écriture de $y$ par l'effet $A$ change
le \emph{membre} de $y$ que lit le dep de l'effet $B$ est une propriété d'une paire
d'arêtes, que ni une arête ni le kill ne portent ; seul le bras self-churn lit les membres
(\code{can_retrigger}, bras membre de la preuve). Un seul cycle simple est reconstruit par
SCC, et un cycle \code{May} qui recoupe une SCC tout-Must n'est pas rapporté
(\cref{sec:churn-tarjan}).

\subsection{Dette documentaire}

Le code fait foi ; trois textes sont en retard sur lui. Les \og Consequences\fg{} d'\adr{42}
disent encore que la condition d'écrivain unique est inchangée, alors que le dernier
paragraphe et le \S6 amendé la déclarent remplacée. Le paragraphe de kill de
\file{docs/relations.md} (\og under every other effect write of the slot in the
component\fg) date de l'amendement \issue{154} : il ignore les sites de rendu et de memo et
l'exclusion des revivers convergents. Enfin le commentaire d'en-tête du bras self-churn
(\src{src/rules/impls/infinite_loop.rs}{393}{394}) dit de l'\Info{} que \og cross-effect
cycles are not analyzed\fg, ce qui est faux depuis \adr{18} : l'\Info{} n'est plus émise que
lorsqu'aucun cycle ne couvre l'écriture. Ordre de confiance : le code au commit
\snapshotCommit, puis les doc-comments de module de \file{churn.rs} et \file{guards.rs}
(réécrits dans ce commit), puis le dernier amendement d'\adr{42} \S6.

% ---------------------------------------------------------------------------
\begin{resume}
\begin{itemize}
  \item Le \textbf{graphe de churn} (\cref{def:churn-graph}) est un pli sur
        \relation{slot_writers} et \relation{effect_triggers} : nœuds = slots qualifiés,
        arête $x \to y$ = \og un changement de $x$ relance un effet qui stocke une référence
        fraîche dans $y$\fg. \code{Must} = dep exact $\wedge$ \code{Fresh} $\wedge$ bloc
        synchrone sur tous les chemins ; \code{May} sinon ; la colonne \code{phase} décide
        (\cref{tab:churn-phases}), \code{Handler} ne porte rien. Un effet sans deps porte une
        auto-arête ; \code{self_slot} isole le bras self-churn (\adr{20} item 2).
  \item \textbf{Tarjan en deux passes} : sous-graphe \code{Must} (cycles \code{all_must},
        \Error{} via \cref{lem:churn-must-cycle}) puis graphe entier (\Warning), en sautant
        les régions couvertes ; un cycle simple par SCC.
  \item Le \textbf{kill de convergence} retire une arête dont l'écriture tire au plus une fois.
        \adr{18} l'autorisait sous condition d'écrivain unique (FP \issue{39}) ; \adr{42} \S6
        le remplace par \code{converges_under_all_writes} : gardes mortes sous
        l'\emph{ensemble tueur} (écriture propre et écritures synchrones co-exécutées sur la
        chaîne) et sous chaque \emph{reviver vivant}, pris avec ses faits
        \emph{invariants} (\cref{def:churn-kill-set,def:churn-invariance}).
  \item Un \textbf{site} (\cref{def:churn-site}) est toute ligne non-handler d'un corps de
        rendu, d'effet non mount-only ou de memo, plus les lignes étrangères d'un effet
        mount-only d'un enfant que la boucle remonte (\code{stays_mounted}). Un slot
        \code{foreign} n'est jamais tué.
  \item Les sites convergents forment un \textbf{plus petit point fixe}
        (\cref{thm:churn-lfp}) : un reviver prouvé n'est plus vivant. Le plus grand point
        fixe lirait deux sites qui se raniment comme convergents (\cref{thm:churn-lfp-sound},
        \cref{ex:churn-mutual}). L'arête est ensuite jugée sur la \emph{partie référence} de
        sa valeur (\code{reference_part}).
  \item \textbf{Cross-component} : lignes \code{owner = Some}, deps \code{Versioned} donc
        jamais \code{Must}, \code{cross_component} plafonné à \Warning{} ; \code{props_hold}
        interdit de lire des faits sur les props d'un composant où une boucle peut entrer.
  \item \textbf{Frontière} : le graphe ne frappe rien ; \code{must_effect_cycle} et
        \code{must_on_all_paths} re-dérivent la preuve ; l'ancre \code{churn_cycles}
        (\adr{29}) projette le graphe par composant.
  \item Fichiers : \file{src/engine/churn.rs} (\code{ChurnGraph::build}, \code{build_edges},
        \code{find_cycles}, \code{stays_mounted}, \code{can_retrigger}),
        \file{src/engine/guards.rs} (\code{WriteSite}, \code{converges_under_all_writes},
        \code{Invariance}, \code{contradicts}), \file{src/engine/written.rs}
        (\code{reference_part}), \file{src/rules/helpers/cycles.rs} (\code{CycleRow}),
        \file{tests/effect_cycles.rs}.
\end{itemize}
Le \cref{chap:inter-composants} explique d'où viennent les lignes étrangères et les valeurs
\code{Versioned} des props que ce chapitre a supposées ; le \cref{chap:infinite-loop}
assemble les quatre bras de \regle{infinite-loop}.
\end{resume}

\section{Exercices}
\label{sec:churn-exercices}

\begin{exercice}{Prédire le graphe}{churn-predire}
Sans lancer l'analyseur, donner les arêtes (extrémités, force, \code{self_slot},
\code{no_deps}) et les cycles du programme suivant, puis la sévérité attendue.
\begin{lstlisting}[language=TypeScript]
export function Q({ p }: { p: boolean }) {
  const [a, setA] = useState({});
  const [b, setB] = useState({});
  useEffect(() => { if (p) setB({ a }); else setB({}); }, [a]);
  useEffect(() => { setA({ ...a, b }); }, [b, a]);
  return null;
}
\end{lstlisting}
\end{exercice}
\begin{solution}
$a \to b$ : deux écritures \code{Fresh} de \code{b} dans deux blocs dont l'union est sur tous
les chemins, dep exact : \code{Must}. $b \to a$ : \code{Must} (écriture à l'entrée, dep
exact). $a \to a$ : dep \code{a} exact et écriture de \code{a} : \code{Must},
\code{self_slot = true}, et \code{write_can_retrigger} répond \code{true} (l'argument n'est
pas un updater fonctionnel). Le cycle $a \to b \to a$ est tout-Must, mono-composant :
\Error{} sur les deux effets par le bras graphe ; l'arête \code{self_slot} donne en plus une
\Error{} du bras self-churn sur le second effet.
\end{solution}

\begin{exercice}{Un reviver stable}{churn-reviver-stable}
Dans l'\cref{ex:churn-multi-writer}, remplacer \code{setB(null)} par \code{setB(CONST)}
où \code{CONST} est un objet de module. L'arête $a \to b$ est-elle tuée ? Quelle ligne de
\code{converges_under_all_writes} décide ?
\end{exercice}
\begin{solution}
Oui, elle est tuée. Le reviver est vivant (non gardé, non convergent) ; mais sa valeur est
une référence \code{Stable}, truthy, et la garde \og \code{b} falsy\fg{} devient bien $\Bot$
sous cette réécriture : le second \code{dead_once_written}
(\srcline{src/engine/guards.rs}{253}) réussit. Observé sur \file{/tmp/ch15/const_reviver.tsx}
(\code{const CONST = { src: 'const' }} au niveau du module) : plus aucun cycle, seule
reste l'\Info{} \og freshly recreates state \code{`a`} outside its deps\fg{} sur l'effet
qui écrit \code{a}. Piège inversé : l'exercice teste que le lecteur
n'applique pas la condition d'écrivain unique d'\adr{18}, qui aurait gardé l'arête ; la
preuve multi-site l'enlève, à raison, puisque \code{CONST} ne remet jamais \code{b} à vide.
\end{solution}

\begin{exercice}{Plus grand point fixe}{churn-gfp}
Dérouler l'itération \og tous convergents, puis retirer\fg{} sur \file{mutual.tsx}
(\code{if (cond) setS(null); if (!s) setS({...})} dans un seul effet) et sur
\file{null_updater.tsx}. Sur lequel des deux le plus grand point fixe diffère-t-il du plus
petit ? Pourquoi ?
\end{exercice}
\begin{solution}
Sur \file{null_updater.tsx} seulement : chacun des deux sites y meurt sous sa propre
écriture et n'est réfuté que par l'autre, donc $\mathrm{gfp} = \{A, B\}$ et
$\lfp = \emptyset$. Dans \file{mutual.tsx}, la garde \code{cond} du site
\code{setS(null)} est un prop qui tient : elle ne meurt sous aucune écriture de \code{s},
le site est réfuté quel que soit l'état de l'autre, puis il ranime l'autre ; les deux
points fixes sont vides.
\end{solution}

\begin{exercice}{Pourquoi jamais Must}{churn-jamais-must}
Expliquer, à partir de l'\cref{lst:churn-dep-edges}, pourquoi l'arête de
l'\cref{ex:churn-cross} ne peut pas être \code{Must} ni \code{self_slot}, puis proposer ce
qu'il faudrait au moteur pour prouver un must-rerun à travers un prop.
\end{exercice}
\begin{solution}
Le dep \code{value} est un prop : sa valeur est \code{Versioned({(Parent, 0)})}, la ligne
\relation{effect_triggers} est \code{exact = false}, le slot entre donc dans
\code{f.versioned} et la seule poussée possible est \code{push(x, EdgeStrength::May, ...)}
(\srcline{src/engine/churn.rs}{592}). \code{self_slot} exige \code{x.0 == f.comp}
(\srcline{src/engine/churn.rs}{589}) : le nœud appartient à \code{Parent}, l'effet à
\code{Child}. Pour un must-rerun, il faudrait savoir que le prop \emph{est} le slot entier
du parent, et non une valeur seulement versionnée par lui : la voie qu'\adr{18} nomme sans
la prendre, une provenance de prop entier au site d'appel JSX.
\end{solution}

\begin{exercice}{Un enfant dans une liste}{churn-map}
Modifier \file{stays_mounted.tsx} pour rendre \code{<Child/>} depuis
\code{items.map(i => <Child key={i} onReady={setS}/>)} avec \code{items} un prop. Que
répond \code{stays_mounted} et pourquoi ? Est-ce un FP ou un FN potentiel ?
\end{exercice}
\begin{solution}
\code{false} : l'élément est dans une closure (\code{ok &= !closure},
\srcline{src/engine/churn.rs}{405}) --- la boucle peut ajouter un élément et monter un
nouvel enfant. La ligne étrangère devient un site, \code{(Parent, 0)} est \code{foreign},
l'arête \code{self_slot} survit : \Warning. Côté fail-closed de \issue{162}, un FP possible,
jamais un FN. Observé (\file{/tmp/ch15/map_child.tsx}) : le \Warning{} \og this effect may
store a fresh reference into state \code{`s`} which its deps react to\fg{}, là où
\file{stays_mounted.tsx} est \og verified\fg.
\end{solution}

\begin{exercice}{Mirror state}{churn-mirror}
(Avancé.) Proposer une approche de \issue{157} qui rende \code{setS(derived)} avec
\code{derived = useMemo(() => ({...}), [s])} visible du graphe sans transformer
\code{setLocal(propValue)} en \Warning. Indication : distinguer \code{Versioned} par un
slot \emph{du même composant} écrit par l'effet, et s'appuyer sur la preuve de convergence
plutôt que sur la fraîcheur.
\end{exercice}
